\documentclass[11pt]{amsart}
\usepackage{amsmath,amssymb,amsthm,mathtools}
\usepackage[margin=1.05in]{geometry}
\usepackage[colorlinks=true,linkcolor=blue,citecolor=blue,urlcolor=blue]{hyperref}
\usepackage{booktabs}

\newtheorem{theorem}{Theorem}[section]
\newtheorem{proposition}[theorem]{Proposition}
\newtheorem{lemma}[theorem]{Lemma}
\newtheorem{corollary}[theorem]{Corollary}
\newtheorem{conjecture}[theorem]{Conjecture}
\theoremstyle{definition}
\newtheorem{definition}[theorem]{Definition}
\newtheorem{remark}[theorem]{Remark}
\numberwithin{equation}{section}

\newcommand{\R}{\mathbb{R}}
\newcommand{\Z}{\mathbb{Z}}
\newcommand{\N}{\mathbb{N}}
\newcommand{\Q}{\mathbb{Q}}
\newcommand{\E}{\mathbb{E}}
\newcommand{\Prob}{\mathbb{P}}
\newcommand{\cH}{\mathcal{H}}
\newcommand{\cC}{\mathcal{C}}
\newcommand{\cF}{\mathcal{F}}
\newcommand{\cI}{\mathcal{I}}
\newcommand{\cB}{\mathcal{B}}
\newcommand{\cW}{\mathcal{W}}
\newcommand{\fp}{\mathfrak{p}}
\newcommand{\fr}[1]{\{#1\}}
\newcommand{\eps}{\varepsilon}
\newcommand{\ind}{\mathbf{1}}
\newcommand{\Prim}{\mathrm{Prim}}
\DeclareMathOperator{\osc}{osc}

\title[Kesten's limit law at a fixed rational starting point]{The discrepancy of $\{n\alpha+\beta\}$ at a fixed rational starting point:\\ Kesten's question and Cauchy laws with arithmetic scales}
\author{Verifiable Mathematical Reasoning Corpus}
\date{October 2026}

\begin{document}
\begin{abstract}
For $\alpha,\beta\in\R$ and $0<r<1$ let $D_N(\alpha,\beta;r)=\#\{1\le n\le N:\fr{n\alpha+\beta}<r\}-Nr$ be the local discrepancy of the
sequence $(\fr{n\alpha+\beta})_{n\ge1}$. Kesten proved in 1960 that $D_N/\log N$ converges in distribution to a Cauchy law when $(\alpha,\beta)$
is uniformly distributed on the unit square, and raised the question whether such a limit law holds for a \emph{fixed} starting point
$\beta$; this question is recorded as open in the recent literature. We prove that for every rational starting point $\beta$ and every rational
$r\in(0,1)$ the normalised discrepancy $D_N(\alpha,\beta;r)/\log N$, with $\alpha$ uniformly distributed on $(0,1)$, converges to a centred
Cauchy law, and we determine its scale: it is $\sigma(\beta,r)=\frac6\pi\,\E\,F(\fr{z_1c},\fr{z_2c})$, where $z_1=\fr{-\beta}$, $z_2=\fr{r-\beta}$,
$F(x,y)=\int_0^1|\tau(t,y)-\tau(t,x)|\,dt$ with the Green function $\tau(t,z)=\min(t,z)(1-\max(t,z))$, and $c$ is the second coordinate of a
uniformly distributed primitive residue vector modulo the common denominator of $z_1,z_2$. In particular $\sigma(0,r)=1/(2\pi)$ for every
rational $r$, whereas Kesten's scale for a random starting point is $1/(3\pi)$; other values include $1/(4\pi)$, $3/(8\pi)$ and, for some
pairs such as $(\beta,r)=(\frac13,\frac12)$, Kesten's value $1/(3\pi)$ itself. So a Cauchy law survives at every fixed rational starting point,
but in general with a scale different from Kesten's.
The proof combines three ingredients: an exact reciprocity formula expressing $D_N$ as an alternating Birkhoff sum along the orbit of
$(\alpha,\fr{N\alpha},z)$ under the skew product $(x,y,z)\mapsto(\fr{1/x},\fr{-y/x},\fr{-z/x})$; the conditional law of the fibre variable
$\fr{N_k\alpha_k}$ given a filtration by renormalisation cells, obtained from an exactly closed two-piece transfer; and, for the interval
endpoints, a finite extension of a Markov refinement of the Gauss map whose exponential mixing we prove by a coupling argument. The value
$1/(2\pi)$ at $\beta=0$ comes from the identity $\sum B_2(c/Q)=0$, the sum running over primitive residue vectors $(b,c)$ modulo $Q$.
\end{abstract}

\maketitle
\section{Introduction}

For real $t$ let $\fr{t}=t-\lfloor t\rfloor$. For $\alpha,\beta\in\R$, $0<r<1$ and $N\ge1$ consider the local discrepancy
\[
D_N(\alpha,\beta;r)=\#\{1\le n\le N:\ \fr{n\alpha+\beta}<r\}-Nr ,\qquad D_N(\alpha;r):=D_N(\alpha,0;r).
\]
For fixed irrational $\alpha$ the sequence $D_N$ is unbounded unless $r\in\Z+\alpha\Z$ (Hecke, Ostrowski, Kesten), and it oscillates on the
scale $\log N$. Kesten \cite{Ke60} discovered that this oscillation has a universal statistical description: if $(\alpha,\beta)$ is uniformly
distributed on $[0,1]^2$, then $D_N(\alpha,\beta;r)/\log N$ converges in distribution to a centred Cauchy law. Kesten expressed the scale
by a complicated formula which appeared to depend on the arithmetic nature of $r$; Borda \cite{Bo25} showed that it equals $1/(3\pi)$ for all
$r\in(0,1)$, and that the analogous law for the sawtooth function $\fr{x}-\frac12$ has scale $1/(4\pi)$. In Kesten's theorem the
randomness comes from \emph{two} parameters. Whether a limit law persists for a \emph{fixed} starting point $\beta$, so that all the
randomness has to come from $\alpha$, ``was already raised by Kesten, and remains open'' \cite{Bo25}; it is also recorded as open in
\cite{Bo23a}.
For the sawtooth function and $\beta=0$ this is Erd\H{o}s Problem \#1002, for which proofs of a Cauchy law with scale $1/(2\pi)$ have
recently been claimed by Kwon and by Wang \cite{Kw26}, and given independently in \cite{E1002}. For the indicator of an interval, i.e.\
for the discrepancy itself, the limit law at a fixed starting point has not been determined. Borda \cite{Bo25} determined, for $\beta=0$ and rational $r$, the
limit laws of $\max_{N<M}D_N$ and $\min_{N<M}D_N$, and asked whether this extends to irrational $r$; these maximal sums only
involve $\alpha$, whereas $D_N$ at a deterministic time $N$ depends on the interplay of $N$ with the continued fraction of $\alpha$. If the time is
randomised as well ($N$ uniform in $\{1,\dots,M\}$), limit laws at a fixed starting point are available for the sawtooth function through the
work of Dolgopyat and Sarig \cite{DS20} (see the discussion in \cite{Bo23a}); the present paper concerns a deterministic time $N$.

\subsection*{The result}
For an integer $Q\ge1$ let $\Prim_Q=\{(b,c)\in(\Z/Q\Z)^2:\gcd(b,c,Q)=1\}$ be the set of primitive residue vectors, and let $c_Q$ denote the
second coordinate of a uniformly distributed element of $\Prim_Q$. For $t,z\in[0,1]$ let
\[
\tau(t,z)=\min(t,z)\,\big(1-\max(t,z)\big)
\]
(the Green function of $-d^2/dt^2$ on $[0,1]$ with Dirichlet boundary conditions), and for $x,y\in[0,1)$ let
\[
F(x,y)=\int_0^1\big|\tau(t,y)-\tau(t,x)\big|\,dt=\frac{|y-x|\,\big(\min(x,y)^2+(1-\max(x,y))^2\big)}{2\,(1-|y-x|)} .
\]
(The closed form is elementary: for $x<y$ the integrand is $t(y-x)$ on $[0,x]$, $(y-x)(1-t)$ on $[y,1]$ and $|t(1-y+x)-x|$ on $[x,y]$.)

\begin{theorem}\label{thm:main}
Let $\beta\in\Q$ and $r\in\Q\cap(0,1)$. Put $z_1=\fr{-\beta}$, $z_2=\fr{r-\beta}$, let $Q\ge2$ be the least common denominator of $z_1$ and
$z_2$, and
\[
\sigma(\beta,r)=\frac6\pi\,\E\,F\big(\fr{z_1c_Q},\fr{z_2c_Q}\big)\ >0 .
\]
Then for every $c\in\R$
\[
\lim_{N\to\infty}\Big|\Big\{\alpha\in(0,1):\ \frac{D_N(\alpha,\beta;r)}{\log N}\le c\Big\}\Big|=\frac12+\frac1\pi\arctan\frac{c}{\sigma(\beta,r)} ;
\]
that is, if $\alpha$ is uniformly distributed on $(0,1)$ then $D_N(\alpha,\beta;r)/\log N$ converges in distribution to the centred Cauchy law
with scale $\sigma(\beta,r)$ (characteristic function $e^{-\sigma(\beta,r)|\theta|}$).
\end{theorem}

The quantity $\sigma(\beta,r)$ is a finite average: $\Prob(c_Q=c)=\varphi(d)/(d\,Q\prod_{p\mid Q}(1-p^{-2}))$ with $d=\gcd(c,Q)$, so $\pi\sigma(\beta,r)$
is a rational number which can be computed exactly. Its positivity is clear, since $c_Q=1$ has positive probability and
$\fr{z_1}\neq\fr{z_2}$.

\begin{corollary}\label{cor:zero}
For every rational $r\in(0,1)$, $D_N(\alpha;r)/\log N$ converges in distribution to the Cauchy law with scale $1/(2\pi)$, i.e.
\[
\lim_{N\to\infty}\Big|\Big\{\alpha\in(0,1):\ \#\{n\le N:\fr{n\alpha}<r\}-Nr\le c\log N\Big\}\Big|=\frac12+\frac1\pi\arctan(2\pi c)\qquad(c\in\R).
\]
\end{corollary}

\begin{proof}
For $\beta=0$ we have $z_1=0$, $z_2=r$, and $Q$ is the denominator of $r$. Since $F(0,y)=\int_0^1\tau(t,y)\,dt=y(1-y)/2$, and
$\E\big[\fr{rc_Q}(1-\fr{rc_Q})\big]=\frac16$ by Lemma~\ref{lem:moebius} below (applied with $P=rQ$, which is coprime to $Q$), we get
$\sigma(0,r)=\frac6\pi\cdot\frac1{12}=\frac1{2\pi}$.
\end{proof}

The scale genuinely depends on the starting point. Exact values of $\pi\sigma(\beta,r)$ for some pairs are:
\begin{center}\small
\begin{tabular}{lcccccccc}\toprule
$(\beta,r)$ & $(0,r)$ & $(\frac14,\frac12)$ & $(\frac15,\frac25)$ & $(\frac23,\frac13)$ & $(\frac17,\frac37)$ & $(\frac13,\frac12)$ & $(\frac12,\frac13)$ & $(\frac7{10},\frac3{10})$\\\midrule
$\pi\sigma(\beta,r)$ & $\frac12$ & $\frac14$ & $\frac14$ & $\frac14$ & $\frac3{10}$ & $\frac13$ & $\frac38$ & $\frac{499}{1512}$\\\bottomrule
\end{tabular}
\end{center}
Thus the answer to Kesten's question, for rational parameters, is: \emph{a Cauchy limit law holds at every fixed rational starting point, but
its scale is an arithmetic function of $(\beta,r)$ which in general differs from Kesten's value $1/(3\pi)$}. For $\beta=0$ it is $1/(2\pi)$ for every
$r$, so Kesten's limit law does not hold at the starting point $0$; for some other pairs, for instance $(\frac13,\frac12)$, $(\frac17,\frac12)$ or
$(\frac14,\frac34)$, the scale happens to equal $1/(3\pi)$. Among all pairs with denominators at most $12$ the values of $\pi\sigma(\beta,r)$ lie in
$[\frac14,\frac12]$. The values in the table are confirmed by Monte Carlo simulation in Section~\ref{sec:numerics}.

\begin{remark}[Consistency with Kesten's theorem]\label{rem:kesten}
For a uniformly distributed starting point the pair of endpoints becomes $(Z,\fr{Z+\delta})$ with $Z$ uniform and $\delta=\fr{rc_Q}$, and a
direct computation gives $\int_0^1F(z,\fr{z+\delta})\,dz=\delta(1-\delta)/3$. Since $\E[\delta(1-\delta)]=\frac16$ (Lemma~\ref{lem:moebius}),
the same recipe produces the scale $\frac6\pi\cdot\frac1{18}=\frac1{3\pi}$, which is Kesten's constant as computed by Borda. Likewise, if
$(z_1c,z_2c)$ is replaced by a pair of independent uniform variables, $\E F=\frac1{18}$.
\end{remark}

\begin{remark}[Irrational parameters]\label{rem:irr}
Our proof uses the rationality of $\beta$ and $r$ in one place only: it makes the orbits of the interval endpoints under the renormalisation
a \emph{finite} extension of a Markov refinement of the Gauss map (Section~\ref{sec:orbits}). Everything else applies verbatim to arbitrary
parameters. The numerical evidence of Section~\ref{sec:numerics} supports the following statement, which we do not prove.
\end{remark}

\begin{conjecture}\label{conj:irr}
For all $\beta\in\R$ and $r\in(0,1)$, $D_N(\alpha,\beta;r)/\log N$ converges to a centred Cauchy law. Its scale is $\frac6\pi\E F(X_1,X_2)$,
where $(X_1,X_2)$ is the limit in distribution of $(\fr{z_1q_K},\fr{z_2q_K})$, $q_K$ being the denominator of the $K$-th convergent of $\alpha$,
with $\alpha$ uniform and $K$ uniform in $\{1,\dots,k\}$, $k\to\infty$. In particular the scale is $1/(2\pi)$ for $\beta\in\Z$ and every $r$,
and $1/(3\pi)$ whenever $1,\beta,r$ are linearly independent over $\Q$.
\end{conjecture}

\subsection*{The method}
Our starting point is an exact reciprocity formula for the local discrepancy (Proposition~\ref{prop:expansion}): with $N_0=N$,
$N_{k+1}=\lfloor N_k\alpha_k\rfloor$, $y_k=\fr{N_k\alpha_k}$ and $z_{k+1}=\fr{-z_k/\alpha_k}$,
\[
D_N(\alpha;z)=\sum_{k=0}^{J-1}(-1)^k\,\ell(\alpha_k,y_k,z_k),\qquad \ell(x,y,z)=\frac{\tau(y,z)}{x}+\kappa(x,y,z),\quad|\kappa|\le1,
\]
where $\alpha_k$ are the Gauss iterates of $\alpha$ and $J\approx\frac{12\log2}{\pi^2}\log N$. Thus $D_N$ is an alternating Birkhoff sum along an
orbit of the skew product $(x,y,z)\mapsto(\fr{1/x},\fr{-y/x},\fr{-z/x})$, and $D_N(\alpha,\beta;r)=D_N(\alpha;z_2)-D_N(\alpha;z_1)$ is a difference
of two such sums driven by the same $(\alpha_k,y_k)$. The large terms are $a_{k+1}(\tau(y_k,z_{2,k})-\tau(y_k,z_{1,k}))$ with $a_{k+1}$ the
partial quotients. The proof has to control three kinds of randomness.

\emph{(i) The digits.} Large partial quotients produce, after normalisation by $\log N$, a Poisson process of jumps with intensity
$\frac{c_0}{\log2}\,dx/x^2$, $c_0=\frac{12\log2}{\pi^2}$; this is where the Cauchy law comes from.

\emph{(ii) The fibre variable $y_k$}, which encodes the position of $N$ relative to the continued fraction of $\alpha$. Its conditional law
given the past is not uniform. We use the description of this conditional law obtained in \cite{E1002}: on a filtration by
``renormalisation cells'', the fibre map transforms piecewise constant densities with pieces $[0,1-x)$, $[1-x,1)$ into densities of the
same kind, the mass of the upper piece following a scalar affine recursion with a uniform two-step contraction. We reproduce the
statements and proofs (Sections~\ref{sec:cells}--\ref{sec:transfer}) to keep the paper self-contained.

\emph{(iii) The endpoint variables $z_{i,k}$}, which do not depend on $N$. For a rational endpoint $z=P/Q$ we show that
$z_k=\fr{u_k+v_k\alpha_k}$, where $u_k=\fr{(-1)^kPq_k/Q}$ and $v_k$ is one of finitely many rationals (Lemma~\ref{lem:orbit}). The pair
$(u_k,v_k)$ is driven by the digits \emph{and} by the position of $\alpha_k$ relative to finitely many rational thresholds; refining the
Gauss partition at these thresholds gives a Markov partition for which $(u_k,v_k)$ is adapted. We prove exponential mixing of the resulting
finite extension by a coupling argument which uses large partial quotients to reset the integer parts of $v_k$ and the fact that the
matrices $\left(\begin{smallmatrix}0&-1\\-1&-c\end{smallmatrix}\right)$, $c\in\Z/Q$, generate $\mathrm{SL}_2(\Z/Q\Z)$ up to sign of the
determinant (Section~\ref{sec:orbits}). The dynamics of $(u_k,v_k)$ does not depend on the parity of $k$, so even and odd levels have the
same limiting statistics; this is what centres the limit law. The stationary law of $u_k$ is that of $\fr{Pc_Q/Q}$, and the elementary
identity
\[
\sum_{(b,c)\in\Prim_Q}B_2\Big(\Big\{\frac cQ\Big\}\Big)=0\qquad(Q\ge2),\qquad B_2(x)=x^2-x+\tfrac16,
\]
shows that the scale at $\beta=0$ does not depend on $r$.

The rest of the argument follows \cite{E1002}: the bounded part of the sum is controlled by a martingale decomposition along the cell
filtration, now with discontinuous summands, and the residual drift by the mixing of the extended chain; the large partial quotients are
handled by a conditional characteristic function argument and a decoupling of the near, transition and far future.

\subsection*{Organisation}
Section~\ref{sec:renorm} proves the reciprocity expansion. Section~\ref{sec:orbits} studies the endpoint orbits for rational endpoints and
proves the mixing of the extended chain. Sections~\ref{sec:cells}--\ref{sec:transfer} describe the conditional law of $y_k$.
Section~\ref{sec:cf} collects facts on continued fractions. Sections~\ref{sec:decomp}--\ref{sec:heavy} carry out the decomposition, the
bounded part and the large partial quotients, and Section~\ref{sec:end} completes the proof. Section~\ref{sec:numerics} reports exact and
Monte Carlo checks.

\subsection*{Notation}
$\lambda$ is Lebesgue measure on $(0,1)$, $\alpha$ is $\lambda$-distributed, and $\Prob,\E$ refer to $\lambda$. For $x\in\R$, $\|x\|$ is the
distance to the nearest integer; for $u,u'\in\R/\Z$ we write $\|u-u'\|$ for the circle distance. Constants $C,C_0,C_1,\dots$ are absolute
unless their dependence on $Q$ is indicated; $c_0=\frac{12\log2}{\pi^2}$.
\section{The reciprocity expansion}\label{sec:renorm}

Throughout, $\alpha\in(0,1)$ is irrational, $\alpha=[0;a_1,a_2,\dots]$, $G(x)=\fr{1/x}$ is the Gauss map,
$\alpha_k=G^k\alpha$ (so $a_{k+1}=\lfloor 1/\alpha_k\rfloor$ and $1/\alpha_k=a_{k+1}+\alpha_{k+1}$), and
$p_k/q_k$ are the convergents, $q_{-1}=0$, $q_0=1$, $q_{k+1}=a_{k+1}q_k+q_{k-1}$, and similarly for $p_k$ with $p_{-1}=1$, $p_0=0$.
We use the classical facts
\begin{equation}\label{eq:theta}
|q_k\alpha-p_k|=\alpha_0\alpha_1\cdots\alpha_k=:|\theta_k|,\qquad \frac{1}{2q_{k+1}}<|\theta_k|<\frac1{q_{k+1}},\qquad
\alpha_k\alpha_{k+1}<\frac12 ,
\end{equation}
and, on the cylinder $\{a_1,\dots,a_k\text{ fixed}\}$, $\alpha=\frac{p_k+p_{k-1}\alpha_k}{q_k+q_{k-1}\alpha_k}$.

\subsection*{The fibre variable of $N$}
Fix $N\ge1$ and define recursively $N_0=N$ and, for $k\ge0$,
\[
Y_k=N_k\alpha_k,\qquad N_{k+1}=\lfloor Y_k\rfloor,\qquad y_k=\fr{Y_k},\qquad J=\min\{k:\ N_k=0\}.
\]
Since $N_{k+1}\le N_k\alpha_k<N_k$ as long as $N_k\ge1$, $J$ is finite; $y_k=0$ for $k\ge J$.

\begin{lemma}\label{lem:yrec}
$y_{k+1}=\fr{-y_k/\alpha_k}$ for all $k\ge0$.
\end{lemma}

\begin{proof}
$N_{k+1}=Y_k-y_k$ gives $Y_{k+1}=N_{k+1}(1/\alpha_k-a_{k+1})=N_k-a_{k+1}N_{k+1}-y_k/\alpha_k$, and $N_k-a_{k+1}N_{k+1}\in\Z$.
\end{proof}

\begin{lemma}[Size of the renormalised times]\label{lem:sizes}
For all $k\ge0$,
\begin{equation}\label{eq:Yk}
Y_k=N|\theta_k|-\sum_{i<k}y_i\,\alpha_{i+1}\cdots\alpha_k,\qquad\text{so}\qquad N|\theta_k|-3<Y_k\le N|\theta_k|.
\end{equation}
Consequently $\frac{N}{2q_{k+1}}-4<N_{k+1}<\frac N{q_{k+1}}$. If $m=\max\{j:q_j\le N\}$, then $m-8\le J\le m+1$.
\end{lemma}

\begin{proof}
$Y_0=N\alpha=N|\theta_0|$ and $Y_k=(Y_{k-1}-y_{k-1})\alpha_k$, which gives \eqref{eq:Yk} by induction. A product of
$t$ consecutive $\alpha_j$ is at most $2^{-\lfloor t/2\rfloor}$ by \eqref{eq:theta}, so
$\sum_{i<k}\alpha_{i+1}\cdots\alpha_k<1+2\sum_{t\ge1}2^{-t}=3$. The bounds on $N_{k+1}=\lfloor Y_k\rfloor$ follow from
\eqref{eq:theta}. If $q_{k+1}>N$ then $N_{k+1}=0$, so $J\le m+1$. If $q_{k+1}\le N/10$ then $N_{k+1}\ge 1$, so $J\ge k+2$;
since $q_{j+2}\ge 2q_j$, at most $8$ indices $j$ satisfy $N/10<q_j\le N$, whence $J\ge m-8$.
\end{proof}

\subsection*{The endpoint orbit and the level function}
For $z\in[0,1)$ put $z_0=z$ and $z_{k+1}=\fr{-z_k/\alpha_k}$. For $x\in(0,1)$ and $y,z\in[0,1)$ write $y'=\fr{-y/x}$, $z'=\fr{-z/x}$ and define
\[
\kappa(x,y,z)=\begin{cases}-(1-z') & \text{if } y\ge z,\\ -(1-y')+\ind\{y'<z'\} & \text{if } y<z,\end{cases}\qquad
\ell(x,y,z)=\frac{\tau(y,z)}{x}+\kappa(x,y,z),
\]
with $\tau(y,z)=\min(y,z)(1-\max(y,z))$ as in the introduction. Note $|\kappa|\le1$, $0\le\tau\le\frac14$, $\tau(y,0)=0$, and $\tau$ is
$1$-Lipschitz in each variable on the circle $\R/\Z$ (it is continuous at $y\in\{0,1\}$ and at $z\in\{0,1\}$, where it vanishes). Write
$D_N(\alpha;z)=D_N(\alpha,0;z)$ and $D_0(\alpha;z)=0$.

\begin{lemma}[Local reciprocity]\label{lem:recip}
Let $z\in(0,1)$ with $z\notin\Z+\alpha\Z$. Then for every $N\ge1$,
\[
D_N(\alpha;z)=\ell\big(\alpha,\fr{N\alpha},z\big)-D_{\lfloor N\alpha\rfloor}\big(G\alpha;\fr{-z/\alpha}\big).
\]
Moreover $\fr{-z/\alpha}\notin\Z+G\alpha\,\Z$.
\end{lemma}

\begin{proof}
Put $M=\lfloor N\alpha\rfloor$, $y=\fr{N\alpha}$, $\alpha_1=G\alpha$, $\rho=\fr{z/\alpha}$ and $z_1=\fr{-z/\alpha}$, so $z_1=1-\rho$ (as
$z/\alpha\notin\Z$), and $y_1=\fr{-y/\alpha}$. Let $C_N(\alpha;z)=\#\{1\le n\le N:\fr{n\alpha}<z\}$. Every $n\in[1,N]$ has $m:=\lfloor n\alpha\rfloor\in[0,M]$,
and $\fr{n\alpha}<z$ iff $m/\alpha\le n<(m+z)/\alpha$. Hence $C_N=\sum_{m=0}^M\#A_m$ with
$A_m=\{n\in\Z\cap[1,N]:\ m/\alpha\le n<(m+z)/\alpha\}$. We use that an interval $[x,x+h)$ with $h>0$ contains exactly
$\lfloor h\rfloor+\ind\{\fr{-x}<\fr{h}\}$ integers. The hypothesis $z\notin\Z+\alpha\Z$ excludes all the equalities between fractional parts that
occur below: $\fr{-m/\alpha}=\rho$ would give $m+z\in\alpha\Z$, and $y_1=z_1$ would give $z-y=z-N\alpha+M\in\alpha\Z$.

\emph{Case $M=0$.} Then $y=N\alpha$ and $C_N=\#\{1\le n\le N: n<z/\alpha\}=\min(N,\lfloor z/\alpha\rfloor)$. If $y\ge z$, then
$C_N=\lfloor z/\alpha\rfloor=z/\alpha-1+z_1$ and $D_N=z(1-y)/\alpha-(1-z_1)=\ell(\alpha,y,z)$. If $y<z$, then $C_N=N$, $D_N=N(1-z)=y(1-z)/\alpha$,
and $y_1=\fr{-N}=0<z_1$, so again $D_N=\ell(\alpha,y,z)$.

\emph{Case $M\ge1$.} For $1\le m\le M-1$ we have $(m+z)/\alpha<M/\alpha\le N$, so $\#A_m=\lfloor z/\alpha\rfloor+\ind\{\fr{-m\alpha_1}<\rho\}$, because
$\fr{-m/\alpha}=\fr{-m\alpha_1}$. As $\fr{-m\alpha_1}=1-\fr{m\alpha_1}$, this indicator equals $1-\ind\{\fr{m\alpha_1}<z_1\}$. Further
$\#A_0=\lceil z/\alpha\rceil-1=\lfloor z/\alpha\rfloor$. Since $\fr{M\alpha_1}=\fr{M/\alpha}=\fr{N-y/\alpha}=y_1$, we obtain
\[
\sum_{m=0}^{M-1}\#A_m=M\lfloor z/\alpha\rfloor+(M-1)-C_{M-1}(\alpha_1;z_1)=M\lfloor z/\alpha\rfloor+M-1-C_M(\alpha_1;z_1)+\ind\{y_1<z_1\}.
\]
For $m=M$ note $N=(M+y)/\alpha$ and $\fr{-M/\alpha}=\fr{y/\alpha}=1-y_1$ (here $y/\alpha\notin\Z$, since $M\ge1$ and $\alpha$ is irrational). If $y\ge z$,
then $(M+z)/\alpha\le N$ and $\#A_M=\lfloor z/\alpha\rfloor+\ind\{1-y_1<\rho\}=\lfloor z/\alpha\rfloor+\ind\{y_1>z_1\}$. If $y<z$, then
$\#A_M=N-\lceil M/\alpha\rceil+1=N-M/\alpha-(1-y_1)+1=y/\alpha+y_1$. Now $C_N=\sum_{m\le M}\#A_m$, $C_M(\alpha_1;z_1)=D_M(\alpha_1;z_1)+Mz_1$ and
$M(\lfloor z/\alpha\rfloor+1-z_1)-Nz=Mz/\alpha-Nz=-yz/\alpha$ give
\[
D_N(\alpha;z)=-D_M(\alpha_1;z_1)-\frac{yz}{\alpha}-1+\ind\{y_1<z_1\}+\#A_M .
\]
If $y\ge z$ the last four terms add up to $\lfloor z/\alpha\rfloor-yz/\alpha=z(1-y)/\alpha-(1-z_1)$; if $y<z$ they add up to
$y(1-z)/\alpha-(1-y_1)+\ind\{y_1<z_1\}$. In both cases this is $\ell(\alpha,y,z)$.

Finally, if $z_1=j+k\alpha_1$ with $j,k\in\Z$, then $-z/\alpha\in\Z+k(1/\alpha-a_1)$ and $z\in\Z+\alpha\Z$, a contradiction.
\end{proof}

\begin{proposition}[Alternating expansion]\label{prop:expansion}
Let $z\in(0,1)\setminus(\Z+\alpha\Z)$, in particular any rational $z\in(0,1)$. For every $N\ge1$,
\begin{equation}\label{eq:expansion}
D_N(\alpha;z)=\sum_{k=0}^{J-1}(-1)^k\ell_k(z),\qquad \ell_k(z):=\ell(\alpha_k,y_k,z_k).
\end{equation}
Moreover $\ell_k(z)=(a_{k+1}+\alpha_{k+1})\tau(y_k,z_k)+\kappa_k(z)$ with $\kappa_k(z)=\kappa(\alpha_k,y_k,z_k)\in[-1,1]$, and $|\ell_k(z)|\le a_{k+1}+1$.
In $\kappa_k(z)$ the quantities $y'$, $z'$ are $y_{k+1}$ and $z_{k+1}$.
\end{proposition}

\begin{proof}
Apply Lemma~\ref{lem:recip} to $(\alpha_k,N_k,z_k)$ for $k=0,\dots,J-1$; the hypothesis propagates by the last statement of the lemma, and the
recursion stops because $D_0=0$. The identification $y'=y_{k+1}$ is Lemma~\ref{lem:yrec}. Finally $|\ell_k|\le(a_{k+1}+1)/4+1\le a_{k+1}+1$.
\end{proof}

Thus $D_N(\alpha;z)$ is the alternating Birkhoff sum of $\ell$ along the orbit of $(\alpha,\fr{N\alpha},z)$ under the skew product
$(x,y,z)\mapsto(\fr{1/x},\fr{-y/x},\fr{-z/x})$, stopped at the absorption time $J$ of the second coordinate. For general rational
starting points we use the following reduction.

\begin{lemma}\label{lem:reduction}
Let $\beta\in\R$, $0<r<1$, $z_1=\fr{-\beta}$, $z_2=\fr{r-\beta}$. Then $D_N(\alpha,\beta;r)=D_N(\alpha;z_2)-D_N(\alpha;z_1)$ for all $N$ and all
$\alpha$.
\end{lemma}

\begin{proof}
$\fr{n\alpha+\beta}<r$ iff $\fr{n\alpha}$ lies in the arc of length $r$ starting at $z_1$. If $z_1<z_2$ this arc is $[z_1,z_2)$ and the count is
$C_N(\alpha;z_2)-C_N(\alpha;z_1)$, while $N(z_2-z_1)=Nr$. If $z_2<z_1$ (in particular if $z_2=0$) the arc is $[z_1,1)\cup[0,z_2)$, the count is
$N-C_N(\alpha;z_1)+C_N(\alpha;z_2)$, and $N-Nz_1+Nz_2-Nr=0$ because $z_2=z_1+r-1$.
\end{proof}

When $\beta,r\in\Q$, Proposition~\ref{prop:expansion} applies to each nonzero endpoint, and with the convention $\ell_k(0):=0$ we obtain for
irrational $\alpha$
\begin{equation}\label{eq:twoexp}
D_N(\alpha,\beta;r)=\sum_{k=0}^{J-1}(-1)^kL_k,\qquad L_k:=\ell_k(z_2)-\ell_k(z_1),\qquad |L_k|\le2(a_{k+1}+1).
\end{equation}
\section{Rational endpoints: a finite extension of a Markov refinement}\label{sec:orbits}

Throughout this section $Q\ge2$ is fixed, and we consider one or two endpoints $z^{(i)}=P_i/Q\in(0,1)$, $P_i\in\{1,\dots,Q-1\}$, with
$\gcd(P_1,P_2,Q)=1$ (with one endpoint, $\gcd(P_1,Q)=1$). In the situation of Theorem~\ref{thm:main} these are the nonzero elements of
$\{z_1,z_2\}$ and $Q$ is their least common denominator. Write $z^{(i)}_k$ for the orbit of $z^{(i)}$ (Section~\ref{sec:renorm}).
Constants in this section may depend on $Q$ and $P_i$.

\subsection{The structure of the endpoint orbits}
Let $w_k=\big((-1)^kq_{k-1},(-1)^kq_k\big)\bmod Q\in\Prim_Q$ (it is primitive since $\gcd(q_{k-1},q_k)=1$). Then $w_0=(0,1)$ and
\begin{equation}\label{eq:wrec}
w_{k+1}=w_kA_{a_{k+1}},\qquad A_a=\begin{pmatrix}0&-1\\-1&-a\end{pmatrix},
\end{equation}
since $w_kA_a=(-1)^k(-q_k,-q_{k-1}-aq_k)$. The matrix $A_a$ depends only on $a\bmod Q$, and $\det A_a=-1$.

\begin{lemma}\label{lem:orbit}
Let $z=P/Q\in(0,1)$ and let $\alpha$ be irrational. There are $v_k\in Q^{-1}\Z$ such that, with $u_k=\fr{P(w_k)_2/Q}=\fr{(-1)^kPq_k/Q}$,
\begin{enumerate}
\item[(i)] $z_k=\fr{u_k+v_k\alpha_k}$ and $z_k\in(0,1)$;
\item[(ii)] $v_k\equiv P(w_k)_1/Q\pmod 1$;
\item[(iii)] $v_0=0$, $v_{k+1}=n_k-u_k$ with $n_k=\lfloor u_k+v_k\alpha_k\rfloor$, and $u_{k+1}=\fr{-u_ka_{k+1}-v_k}$;
\item[(iv)] $-2-2\alpha_{k-1}<v_k\le0$ for $k\ge1$; in particular $-4<v_k\le0$ for all $k$.
\end{enumerate}
\end{lemma}

\begin{proof}
Define $v_k$ by (iii). Then (i)--(ii) hold for $k=0$. Assume them for $k$ and write $z_k=u_k+v_k\alpha_k-n_k$. Since $1/\alpha_k=a_{k+1}+\alpha_{k+1}$,
\[
-\frac{z_k}{\alpha_k}=-(u_k-n_k)(a_{k+1}+\alpha_{k+1})-v_k\equiv\big(-u_ka_{k+1}-v_k\big)+(n_k-u_k)\alpha_{k+1}\pmod1 ,
\]
so $z_{k+1}=\fr{u_{k+1}+v_{k+1}\alpha_{k+1}}$ provided $u_{k+1}\equiv-u_ka_{k+1}-v_k$; and indeed, by \eqref{eq:wrec} and (ii),
$P(w_{k+1})_2/Q=P(-(w_k)_1-a_{k+1}(w_k)_2)/Q\equiv-v_k-a_{k+1}u_k$. Also $v_{k+1}=n_k-u_k\equiv-P(w_k)_2/Q=P(w_{k+1})_1/Q$, which is (ii).
That $z_k\ne0$ follows from $z\notin\Z+\alpha\Z$ and Lemma~\ref{lem:recip}.
For (iv): $v_{k+1}=v_k\alpha_k-z_k\in(v_k\alpha_k-1,v_k\alpha_k]$. Hence $v_k\le0$ for all $k$, $v_1=-z>-1$, and if $v_k>-2-2\alpha_{k-1}$ then
$v_{k+1}>-2\alpha_k-2\alpha_{k-1}\alpha_k-1>-2\alpha_k-2$ by \eqref{eq:theta}.
\end{proof}

We write $\zeta_k$ for the vector of the pairs $(u^{(i)}_k,v^{(i)}_k)$ attached to the endpoints $z^{(i)}$. By Lemma~\ref{lem:orbit}(ii) and the
definition of $u^{(i)}_k$, $(Qv^{(i)}_k,Qu^{(i)}_k)\equiv P_iw_k\pmod Q$. Choosing integers $e_i$ with $\sum_ie_iP_i\equiv1\pmod Q$ (possible since
$\gcd(P_1,P_2,Q)=1$) we get $w_k\equiv\sum_ie_i(Qv^{(i)}_k,Qu^{(i)}_k)$: \emph{the residue vector $w_k$ is a function of $\zeta_k$}.

\subsection{The refined partition}
Let $D_0=4Q$, $E_0=\{j/d:\ 1\le j<d\le D_0\}$ and $B=E_0\cup\{1/a:\ a\ge1\}$. A \emph{piece} is a connected component of $(0,1)\setminus B$; each piece $\fp$
lies in a branch $(\frac1{a+1},\frac1a)$ of $G$, and we put $a(\fp)=a$. For $a\ge D_0$ the branch $(\frac1{a+1},\frac1a)$ contains no point of $E_0$
(a point $j/d<1/D_0$ with $d\le D_0$ would have $j<1$), so it is a single piece; we call these pieces \emph{large}.

\begin{lemma}\label{lem:pieces}
\textup{(a)} For $u\in Q^{-1}\Z\cap[0,1)$ and $v\in Q^{-1}\Z\cap(-4,0]$ the function $x\mapsto\lfloor u+vx\rfloor$ is constant on every piece.
\textup{(b)} $G$ maps every piece $\fp$ increasingly or decreasingly onto an open interval $G(\fp)$ whose endpoints lie in $E_0\cup\{0,1\}$; for a large
piece $G(\fp)=(0,1)$. The set $\mathfrak J=\{G(\fp)\}\cup\{(0,1)\}$ is finite, and every $I\in\mathfrak J$ is, up to a countable set, the union of the
pieces it contains.
\end{lemma}

\begin{proof}
(a) The function jumps only where $u+vx\in\Z$, i.e.\ (for $v\ne0$) at $x=(Qm-Qu)/(Qv)$ with $m\in\Z$; since $0<-Qv<D_0$, such points of $(0,1)$ lie in
$E_0$. (b) The endpoints of a piece in the branch of $a$ are points of $E_0$ or the branch endpoints $\frac1{a+1},\frac1a$, at which $G$ has the
one-sided limits $1$ and $0$. For $j/d\in E_0$ in the interior of the branch, $G(j/d)=\fr{d/j}$ is a fraction with denominator $j<D_0$, hence lies in
$E_0\cup\{0\}$. The set $\mathfrak J$ is finite because $E_0$ is. An open interval with endpoints in $E_0\cup\{0,1\}$ is the union of the pieces it
contains and of the points of $B$ in it.
\end{proof}

For irrational $\alpha$ let $\fp_k$ be the piece containing $\alpha_k$ (note $\alpha_k\notin B$), and let $\cF'_k=\sigma(\fp_0,\dots,\fp_{k-1})$; this
contains $\cF_k=\sigma(a_1,\dots,a_k)$ because $a_{j+1}=a(\fp_j)$. By Lemma~\ref{lem:pieces}(a), $n^{(i)}_k=\lfloor u^{(i)}_k+v^{(i)}_k\alpha_k\rfloor$ is a
function of $(\zeta_k,\fp_k)$, so by Lemma~\ref{lem:orbit}(iii) there are maps $f_\fp$ with
\begin{equation}\label{eq:zetarec}
\zeta_{k+1}=f_{\fp_k}(\zeta_k);\qquad\text{hence $\zeta_k$ is $\cF'_k$-measurable.}
\end{equation}
Let $\omega_k=q_{k-1}/q_k$ ($\omega_0=0$), and $\cI_k=G(\fp_{k-1})$ ($\cI_0=(0,1)$).

\begin{lemma}[Refined cylinders]\label{lem:refcyl}
Let $\fp_0,\dots,\fp_{k-1}$ be pieces and $Z=\{\alpha:\alpha_j\in\fp_j\ (j<k)\}\ne\emptyset$. Then $\alpha\mapsto\alpha_k$ maps $Z$ bijectively onto
$G(\fp_{k-1})$ (onto $(0,1)$ if $k=0$), with inverse $t\mapsto\frac{p_k+p_{k-1}t}{q_k+q_{k-1}t}$, where $p_j,q_j$ are formed with the digits
$a(\fp_0),\dots,a(\fp_{k-1})$. Consequently, conditionally on $\cF'_k$, $\alpha_k$ has density proportional to $(1+\omega_kt)^{-2}$ on $\cI_k$.
\end{lemma}

\begin{proof}
By induction on $k$. If $G^{k-1}$ maps $Z'=\{\alpha_j\in\fp_j,\ j<k-1\}$ onto $G(\fp_{k-2})$, then $\fp_{k-1}$ meets, hence (Lemma~\ref{lem:pieces}(b)) is contained in,
$G(\fp_{k-2})$, so $G^{k-1}$ maps $Z$ onto $\fp_{k-1}$ and $G^k$ maps $Z$ onto $G(\fp_{k-1})$. The inverse is the restriction of the inverse branch of
$G^k$ on the cylinder of the digits, which is the stated M\"obius map; its derivative is $\pm(q_k+q_{k-1}t)^{-2}$.
\end{proof}

\subsection{The extended chain}
For $\omega\in[0,1]$ and $I\in\mathfrak J$ let $K_{\omega,I}$ be the law of the piece containing a point $t\in I$ drawn with density proportional to
$(1+\omega t)^{-2}$. Since $(1+\omega t)^{-2}\in[\frac14,1]$,
\begin{equation}\label{eq:Klower}
K_{\omega,I}(\fp)\ge\frac{|\fp|}{4|I|}\ge\frac{|\fp|}4\qquad(\fp\subset I),\qquad K_{\omega,(0,1)}\big((\tfrac1{a+1},\tfrac1a)\big)\ge\frac1{4a(a+1)} .
\end{equation}
Let $W_k=(\omega_k,\cI_k,\zeta_k)$. By Lemma~\ref{lem:refcyl}, given $\cF'_k$ the piece $\fp_k$ has law $K_{\omega_k,\cI_k}$, and then
$W_{k+1}=(1/(a(\fp_k)+\omega_k),G(\fp_k),f_{\fp_k}(\zeta_k))$. Hence $(W_k)$ is a time-homogeneous Markov chain with respect to $(\cF'_k)$, with
kernel
\[
\mathsf K\big((\omega,I,\zeta),\cdot\big)=\text{law of }\big(1/(a(\fp)+\omega),\,G(\fp),\,f_\fp(\zeta)\big),\qquad \fp\sim K_{\omega,I},
\]
and the conditional law of $(\fp_k,\fp_{k+1},\dots)$ given $\cF'_k$ is the law of the pieces emitted by this chain started from $W_k$.
Let $\cW$ be the (finite) set of pairs $(\cI_k(\alpha),\zeta_k(\alpha))$, $\alpha\in(0,1)\setminus\Q$, $k\ge0$. If $(I,\zeta)=(\cI_k(\alpha),\zeta_k(\alpha))$ and
$\fp\subset I$, then by Lemma~\ref{lem:refcyl} some $\alpha'$ with the same first $k$ pieces has $\alpha'_k\in\fp$, so $(G(\fp),f_\fp(\zeta))\in\cW$.
Thus the chain started from $[0,1]\times\cW$ stays there, and we always consider such starting points.

\begin{lemma}[Generation]\label{lem:gen}
There is an even $n_Q$ such that every $g\in\mathrm{SL}_2(\Z/Q\Z)$ is a product $A_{c_1}\cdots A_{c_{n_Q}}$ with $c_j\in\Z/Q\Z$. Moreover
$\mathrm{SL}_2(\Z/Q\Z)$ acts transitively on $\Prim_Q$ by right multiplication.
\end{lemma}

\begin{proof}
$A_0^2=I$, $A_0A_c=\left(\begin{smallmatrix}1&c\\0&1\end{smallmatrix}\right)$ and $A_cA_0=\left(\begin{smallmatrix}1&0\\c&1\end{smallmatrix}\right)$. The reduction map
$\mathrm{SL}_2(\Z)\to\mathrm{SL}_2(\Z/Q\Z)$ is surjective and $\mathrm{SL}_2(\Z)$ is generated by the two elementary matrices with $c=1$, whose inverses
are the elementary matrices with $c=-1$; so every element of the finite group $\mathrm{SL}_2(\Z/Q\Z)$ is a product of an even number $\le n_Q$ of the
$A_c$, and we pad with $A_0A_0$. For transitivity it suffices to show that every $(b,c)\in\Prim_Q$ is congruent to a primitive vector of $\Z^2$
(which is the second row of a matrix in $\mathrm{SL}_2(\Z)$, and $(0,1)$ times that matrix is the vector). Represent $c$ by an integer $c\ne0$ and let
$t$ be the product of the primes dividing $c$ but not $b$. If a prime $p\mid c$ divides $b$, then $p\nmid Q$ and $p\nmid t$, so $p\nmid b+tQ$; otherwise
$p\mid t$ and $b+tQ\equiv b\not\equiv0\pmod p$. Thus $\gcd(b+tQ,c)=1$.
\end{proof}

\begin{lemma}[Coupling]\label{lem:coupling}
There are constants $C_Q<\infty$ and $\vartheta_Q\in(0,1)$ with the following property. For any two states $W=(\omega,I,\zeta)$ and
$W'=(\omega',I',\zeta')$ with $(I,\zeta),(I',\zeta')\in\cW$ there is a coupling of the chains started at $W$ and $W'$, and a random time $\tau$ with
$\Prob(\tau>n)\le C_Q\vartheta_Q^n$, such that for all $k\ge\tau$ the emitted pieces coincide, $(\cI_k,\zeta_k)=(\cI'_k,\zeta'_k)$, and
$|\omega_k-\omega'_k|\le F_{k-\tau+1}^{-2}$, where $F_j$ are the Fibonacci numbers.
\end{lemma}

\begin{proof}
\emph{Step 1 (reaching the full interval).} Every $I\in\mathfrak J$ contains a cylinder $\{a_1=b_1,\dots,a_m=b_m\}$, and the points of this cylinder
with $a_{m+1}=D_0$ form a non-empty open interval; the pieces of $t,Gt,\dots,G^mt$ for an irrational $t$ in it form an admissible sequence
$\fp^I_1,\dots,\fp^I_{m(I)}$ ending with a large piece. By \eqref{eq:Klower} a chain at any state with interval $I$ emits this sequence with probability
at least $\prod_j|\fp^I_j|/4>0$; afterwards its interval is $(0,1)$, and by \eqref{eq:Klower} it can emit the large piece of digit $D_0$ with
probability $\ge(4D_0(D_0+1))^{-1}$ at each further step. Hence there are $m_1$ and $c_1>0$ such that from any state, with probability
$\ge c_1$, after exactly $m_1$ steps the chain has interval $(0,1)$.

\emph{Step 2 (synchronising $\zeta$).} Consider a chain with interval $(0,1)$ which emits $n_Q+1$ large pieces with digits $b_j\in[D_0,D_0+Q)$. For a
large digit $b$ and $x=\alpha_k\in(\frac1{b+1},\frac1b)$ we have $|v^{(i)}x|<4/D_0=1/Q$, so if $u^{(i)}\ne0$ then $u^{(i)}\in[\frac1Q,1-\frac1Q]$ gives
$\lfloor u^{(i)}+v^{(i)}x\rfloor=0$ and the new $v^{(i)}$ equals $-u^{(i)}$. By Lemma~\ref{lem:gen} and \eqref{eq:wrec}, given the current $w$ we can
choose $b_1,\dots,b_{n_Q}$ (as residues modulo $Q$) such that $w$ is mapped to $(0,1)$; then $u^{(i)}=\fr{P_i/Q}\ne0$ for every endpoint, and after one
more step with $b_{n_Q+1}=D_0$ the new state is $w=(0,1)A_{D_0}=(-1,0)$, $u^{(i)}=0$, $v^{(i)}=-P_i/Q$, independently of the state before the block.
Each such sequence of $n_Q+1$ large pieces is emitted with probability $\ge c_2:=(4(D_0+Q)^2)^{-n_Q-1}$.

\emph{Step 3 (staying together).} Suppose that at some time both chains have the same $(I,\zeta)$. At each subsequent step draw the two pieces from a
maximal coupling of $K_{\omega,I}$ and $K_{\omega',I}$. When they agree, $(I,\zeta)$ stays equal and the common digit $b$ maps $\omega,\omega'$ to
$1/(b+\omega)$, $1/(b+\omega')$. The densities $(1+\omega t)^{-2}/Z_\omega$ satisfy $|\partial_\omega\log((1+\omega t)^{-2}/Z_\omega)|\le4$, so the probability of
disagreement is at most $\min(\frac34,8|\omega-\omega'|)$ (the bound $\frac34$ follows from \eqref{eq:Klower}). After $j$ agreements with digits
$b_1,\dots,b_j$, $\omega$ and $\omega'$ are the values at $\omega_0,\omega_0'$ of $s\mapsto[0;b_j,\dots,b_2,b_1+s]$, whose derivative is at most the inverse
square of the continuant of $b_1,\dots,b_j$, which is $\ge F_{j+1}$. Hence all agreements succeed with probability at least
$c_3=\prod_{j\ge0}\max(\frac14,1-8F_{j+1}^{-2})>0$.

\emph{Renewal.} An \emph{attempt} consists of $m_1+n_Q+1$ steps in which the two chains move independently, followed, if both chains have followed
Steps~1 and~2, by Step~3. A new attempt is started whenever the current one fails (at the end of the independent phase, or at the first disagreement in
Step~3). This is a Markovian coupling with the correct marginals. Each attempt succeeds with conditional probability $\ge\delta=(c_1c_2)^2c_3$ given the
past, and a failed attempt lasts $m_1+n_Q+1+j$ steps, where $j$ is the step of Step~3 at which it fails, with
$\Prob(j\ge j_0)\le\sum_{j\ge j_0}8F_{j+1}^{-2}$. For $\eta>0$ small enough the conditional expectation of $e^{\eta\cdot\text{duration}}$ over a failed
attempt is at most $1-\delta/2$, so the starting time $\tau$ of the successful Step~3 satisfies $\E e^{\eta\tau}<\infty$, uniformly in the initial
states. From $\tau$ on everything agrees, and $|\omega_k-\omega'_k|\le F_{k-\tau+1}^{-2}$ by Step~3.
\end{proof}

\subsection{Consequences}
By a \emph{functional of the chain from time $m$} we mean a random variable $\Phi=\phi(\omega_m,\cI_m,\zeta_m,\fp_m,\fp_{m+1},\dots)$ with $\phi$ Borel; it is
called \emph{$\Lambda$-regular} if $|\phi|\le1$ and $\phi$ is $\Lambda$-Lipschitz in its first argument. $\Phi^{[m']}$ denotes the same $\phi$
evaluated at time $m'$.

\begin{lemma}[Mixing]\label{lem:mixing}
There are constants $C_Q'$ and $\vartheta_Q'\in(0,1)$ such that for every $\Lambda$-regular functional $\Phi$ of the chain from time $m$:
\begin{enumerate}
\item[(a)] $\big|\E[\Phi\mid\cF'_j]-\E\Phi\big|\le C_Q'(1+\Lambda)\vartheta_Q'^{\,m-j}$ almost surely, for $0\le j\le m$;
\item[(b)] the limit $\tilde\mu(\phi)=\lim_{m'\to\infty}\E\Phi^{[m']}$ exists and $|\E\Phi^{[m']}-\tilde\mu(\phi)|\le C_Q'(1+\Lambda)\vartheta_Q'^{\,m'}$.
\end{enumerate}
\end{lemma}

\begin{proof}
(a) By the Markov property $\E[\Phi\mid\cF'_j]=g(W_j)$ with $g(W)=\E_W\Phi^{[m-j]}$, the chain being started at $W$ at time $0$. For two states
coupled as in Lemma~\ref{lem:coupling}, $|g(W)-g(W')|\le2\Prob(\tau>n)+\Lambda\,\E\big[F_{n-\tau+1}^{-2};\tau\le n\big]$ with $n=m-j$, which is
$\le C(1+\Lambda)\vartheta^n$ for some $\vartheta<1$ (split according to $\tau\le n/2$). As $\E\Phi=\E g(W_j)$, the claim follows.
(b) $\E\Phi^{[m'+1]}=\E h(W_1)$ and $\E\Phi^{[m']}=h(W_0)$ with $h(W)=\E_W\Phi^{[m']}$; the same coupling bound shows that $(\E\Phi^{[m']})$ is a Cauchy
sequence with the stated rate.
\end{proof}

Since the kernel does not depend on the parity of the time, \emph{even and odd times have the same limit $\tilde\mu$}; this is the source of
the cancellation of drifts in Sections~\ref{sec:light} and~\ref{sec:heavy}.

\begin{lemma}[The stationary law of $(\omega,w)$]\label{lem:statlaw}
Let $\gamma(dt)=\frac{dt}{(1+t)\log2}$ be the Gauss measure and $\mathrm U_Q$ the uniform law on $\Prim_Q$. For $f:[0,1]\times\Prim_Q\to[-1,1]$
which is $\Lambda$-Lipschitz in the first variable,
$\big|\E f(\omega_k,w_k)-\int f\,d(\gamma\otimes\mathrm U_Q)\big|\le C_Q'(1+\Lambda)\vartheta_Q'^{\,k}$. In particular $c_k:=(w_k)_2$ converges in law to
$c_Q$, and $\tilde\mu\big(g(\omega)h(c)\big)=\int g\,d\gamma\cdot\E h(c_Q)$.
\end{lemma}

\begin{proof}
If $t$ has density $(1+\omega)(1+\omega t)^{-2}$ on $(0,1)$ and $\omega\sim\gamma$, the pair $(t,\omega)$ has density
$(\log2)^{-1}(1+t\omega)^{-2}$, which is invariant under $(t,\omega)\mapsto(Gt,1/(a_1(t)+\omega))$ (the natural extension of $G$; with
$t=1/(a+t')$, $\omega=1/\omega'-a$ one has $1+t\omega=(1+t'\omega')/(\omega'(a+t'))$, and the Jacobian is $(a+t')^{-2}\omega'^{-2}$); its
second marginal is $\gamma$. Hence the digit-driven chain $\omega\mapsto1/(a+\omega)$, $w\mapsto wA_a$, $a\sim$ law of $a_1(t)$, preserves
$\gamma\otimes\mathrm U_Q$ ($w\mapsto wA_a$ is a bijection of $\Prim_Q$). For each $w\in\Prim_Q$ choose $\zeta(w)$ with $((0,1),\zeta(w))\in\cW$
and residue vector $w$ (steer with large digits as in Step~2). Started from $(\omega,(0,1),\zeta(w))$ with $(\omega,w)\sim\gamma\otimes\mathrm U_Q$,
the extended chain emits the pieces of the orbit of a point $t$ with density $(1+\omega)(1+\omega t)^{-2}$: this is the computation of
Lemma~\ref{lem:refcyl} with the prior density of $t$ included, which replaces $q_j$ by $q_j+\omega p_j$ and hence $\omega_k$ by
$(q_{k-1}+\omega p_{k-1})/(q_k+\omega p_k)=[0;a_k,\dots,a_2,a_1+\omega]$, the value produced by the kernel. So its $(\omega_k,w_k)$ is the stationary
digit-driven chain. Coupling it with the actual chain (Lemma~\ref{lem:coupling}) gives the estimate.
\end{proof}

\begin{lemma}\label{lem:moebius}
For $Q\ge2$, $\sum_{(b,c)\in\Prim_Q}B_2\big(\fr{c/Q}\big)=0$, where $B_2(x)=x^2-x+\frac16$. Consequently $\E\big[\fr{Pc_Q/Q}(1-\fr{Pc_Q/Q})\big]=\frac16$ for
every $P$ coprime to $Q$.
\end{lemma}

\begin{proof}
By M\"obius inversion over $d\mid\gcd(b,c,Q)$ and the distribution relation $\sum_{j=0}^{m-1}B_2(j/m)=\frac1{6m}$,
\[
\sum_{(b,c)\in\Prim_Q}B_2(\fr{c/Q})=\sum_{d\mid Q}\mu(d)\,\frac Qd\sum_{j\bmod Q/d}B_2\Big(\frac{j}{Q/d}\Big)=\frac16\sum_{d\mid Q}\mu(d)=0 .
\]
Since $\Prob(c_Q=c)$ depends only on $\gcd(c,Q)$, $Pc_Q$ has the same law as $c_Q$; and $x(1-x)=\frac16-B_2(x)$.
\end{proof}
\section{Cells}\label{sec:cells}

Fix $N$. For $j\ge0$ and irrational $\alpha$ let the \emph{level-$j$ cell} of $\alpha$ be
\[
\cC_j(\alpha)=\{\alpha'\in(0,1)\setminus\mathbb Q:\ a_l(\alpha')=a_l(\alpha),\ N_l(\alpha')=N_l(\alpha)\ (1\le l\le j)\},
\]
so $\cC_0=(0,1)\setminus\mathbb Q$. Let $\cH_j$ be the $\sigma$-algebra generated by the level-$j$ cells. Clearly
$\cH_j\subset\cH_{j+1}$, and $J$ is an $(\cH_j)$-stopping time, since $\{J=j\}=\{N_{j-1}\ge1,N_j=0\}\in\cH_j$.

\begin{lemma}\label{lem:cells}
Every cell is the set of irrational points of an open interval. On a level-$j$ cell $C$ ($j\ge1$) with $N_{j-1}\ge1$ the maps
$\alpha\mapsto\alpha_{j-1}$ and $\alpha\mapsto Y_{j-1}=N_{j-1}\alpha_{j-1}$ are continuous and strictly monotone,
$Y_{j-1}(C)\subset[N_j,N_j+1)$, and hence $y_{j-1}=Y_{j-1}-N_j$ maps $C$ monotonically onto an interval
$R_{j-1}(C)\subset[0,1)$, the \emph{range} of $C$.
\end{lemma}

\begin{proof}
On a cylinder $\{a_1,\dots,a_{j-1}\text{ fixed}\}$ the map $\alpha\mapsto\alpha_{j-1}$ is a M\"obius map,
$\alpha=\frac{p_{j-1}+p_{j-2}\alpha_{j-1}}{q_{j-1}+q_{j-2}\alpha_{j-1}}$. By induction on $j$: inside a level-$(j-1)$ cell
(an interval on which $N_{j-1}$ is constant), $a_j=\lfloor1/\alpha_{j-1}\rfloor$ and $N_j=\lfloor N_{j-1}\alpha_{j-1}\rfloor$
are monotone step functions of the monotone function $\alpha_{j-1}$, so their level sets are intervals.
\end{proof}

\begin{lemma}[Fibre relation]\label{lem:fibre}
Let $C$ be a level-$j$ cell with $N_j\ge1$, and put $M=N_{j-1}-a_jN_j\in\Z$, $\Phi(u)=\frac{N_{j-1}u}{N_j+u}$. On $C$,
\[
\alpha_{j-1}=\frac{N_j+y_{j-1}}{N_{j-1}},\qquad Y_j=M-\Phi(y_{j-1}),\qquad
\frac{1}{\Phi'(y_{j-1})}=\alpha_{j-1}\Big(1+\frac{y_{j-1}}{N_j}\Big).
\]
\end{lemma}

\begin{proof}
$Y_{j-1}=N_j+y_{j-1}=N_{j-1}\alpha_{j-1}$, and $Y_j=N_j(1/\alpha_{j-1}-a_j)=\frac{N_jN_{j-1}}{N_j+y_{j-1}}-a_jN_j=M-\Phi(y_{j-1})$.
Finally $\Phi'(u)=\frac{N_{j-1}N_j}{(N_j+u)^2}$ and $N_j+y_{j-1}=N_{j-1}\alpha_{j-1}$.
\end{proof}

The endpoints of cells are rational. An endpoint $e$ of a level-$j$ cell is a point at which some $a_l$ or some $N_l$
($l\le j$) is discontinuous. We call a point $e$ an \emph{integer point of level $i$} if $a_1,\dots,a_i$ and $N_1,\dots,N_{i-1}$
are constant on a neighbourhood of $e$ and $\lim_{\alpha\to e}Y_{i-1}(\alpha)\in\Z$ (the depth to which the digits are constant near $e$ is
an additional hypothesis in the lemma below). For a function $F$ continuous near
$e$ we write $F(e)$ for its limit at $e$.

\begin{lemma}[Integer points]\label{lem:intpt}
Let $e$ be an integer point of level $i$, $\lim_{\alpha\to e}Y_{i-1}=n+1$, and suppose $a_1,\dots,a_{m+1}$ are constant near
$e$, where $m\ge i$, and that $N_l\ge2$ for all $l\le m$ on one side of $e$ (near $e$). Call $A$ the side of $e$ on which $Y_{i-1}<n+1$ and $B$ the
other side. Then the one-sided limits at $e$ are
\[
\begin{array}{c|ccc}
 & y_{i-1} & y_i & y_l\ (i<l\le m)\\\hline
A & 1 & 1-\alpha_i(e) & 1-\alpha_l(e)\\
B & 0 & 1 & 1-\alpha_l(e)
\end{array}
\]
and $N_l$ is the same on both sides for $l\ne i$, $l\le m+1$ (while $N_i=n$ on $A$, $n+1$ on $B$).
\end{lemma}

\begin{proof}
At $e$, $\alpha_{i-1}(e)=\frac{n+1}{N_{i-1}}$, so $c:=(n+1)\alpha_i(e)=N_{i-1}-a_i(n+1)\in\Z$. On $A$, $N_i=n$, $y_{i-1}\to1$, and
$Y_i=n\alpha_i\to c-\alpha_i(e)$, so $N_{i+1}=c-1$ and $y_i\to1-\alpha_i(e)$. On $B$, $N_i=n+1$, $y_{i-1}\to0$; as
$\alpha_{i-1}$ decreases towards $\alpha_{i-1}(e)$ from above, $\alpha_i$ increases to $\alpha_i(e)$, so $Y_i=(n+1)\alpha_i\uparrow c$, whence
$N_{i+1}=c-1$ and $y_i\to1$. In both cases $-y_i/\alpha_i$ tends to a non-integer
($1-a_{i+1}-\alpha_{i+1}(e)$, resp.\ $-a_{i+1}-\alpha_{i+1}(e)$), so by Proposition~\ref{prop:expansion},
$y_{i+1}\to\fr{-1/\alpha_i(e)}=1-\alpha_{i+1}(e)$ on both sides; inductively $y_l\to1-\alpha_l(e)$ for $i<l\le m$, these limits being
non-integers. Since $N_{i+1}$ agrees on both sides and $Y_l=N_l\alpha_l$ is continuous with non-integer limit for
$l>i$, the $N_l$ agree for $l>i$. In particular, as $N_i\in\{n,n+1\}$ and the other $N_l$ agree, $N_l\ge1$ for $l\le m$ on both sides, so no absorption occurs and the
computation above is legitimate on both sides.
\end{proof}

\begin{definition}\label{def:regular}
Let $1\le i\le k+1$. A level-$i$ cell $C$ is \emph{$k$-regular} if $a_1,\dots,a_{k+2}$ are constant on an open interval containing the
closure of $C$, and $N_l\ge2$ on $C$ for all $l\le k+1$.
\end{definition}

Let $C$ be a $k$-regular level-$i$ cell and $e$ an endpoint of a level-$l$ sub-cell $C'\subset C$, $i\le l\le k+1$. Since $e\in\overline C$,
the digits $a_1,\dots,a_{k+2}$ are constant near $e$, so only some $N_j$ ($j\le l$) is discontinuous at $e$. If $i'$ is the least such
$j$, then $N_1,\dots,N_{i'-1}$ are constant near $e$, $Y_{i'-1}=N_{i'-1}\alpha_{i'-1}$ is continuous near $e$, and $N_{i'}=\lfloor Y_{i'-1}\rfloor$
jumps, so $\lim_{\alpha\to e}Y_{i'-1}\in\Z$: $e$ is an integer point of level $i'\le l$, and since $N_j\ge2$ ($j\le k+1$) on the side of
$e$ inside $C'$, Lemma~\ref{lem:intpt} applies at $e$ with $m=k+1$. Every level-$l$ sub-cell ($i\le l\le k+1$) of a $k$-regular
level-$i$ cell is again $k$-regular.

\begin{proposition}[Three types]\label{prop:types}
Let $C$ be a level-$j$ cell whose endpoints are integer points (of levels $\le j$) satisfying the hypotheses of
Lemma~\ref{lem:intpt} with $m\ge j$, and assume $N_{j-1}\ge2$ on $C$. Then, up to its endpoints, $R_{j-1}(C)$ is one of
\[
\textbf{F}:\ [0,1),\qquad \textbf{P}:\ (1-\alpha_{j-1}(e),1),\qquad \textbf{Q}:\ [0,1-\alpha_{j-1}(e)),
\]
where $e$ is an endpoint of $C$. In particular this holds for every level-$j$ sub-cell of a $k$-regular cell, $j\le k+1$.
\end{proposition}

\begin{proof}
$R_{j-1}(C)$ is the interval between the one-sided limits of $y_{j-1}$ at the two endpoints $e_1,e_2$ of $C$. By
Lemma~\ref{lem:intpt} (applied with the level $i\le j$ of the endpoint, and $l=j-1\ge i-1$) each limit belongs to
$\{0,1,1-\alpha_{j-1}(e_\nu)\}$. Since $Y_{j-1}=N_{j-1}\alpha_{j-1}$ on $C$, $|R_{j-1}(C)|=N_{j-1}|\alpha_{j-1}(e_1)-\alpha_{j-1}(e_2)|>0$. If both
limits were of the third kind we would get $|R_{j-1}(C)|=|R_{j-1}(C)|/N_{j-1}$, impossible for $N_{j-1}\ge2$. If both lie in
$\{0,1\}$ the range is $[0,1)$. Otherwise it is of type \textbf{P} or \textbf{Q}.
\end{proof}

\begin{lemma}[Distortion]\label{lem:distortion}
Let $C$ be a level-$j$ cell, $j\ge1$, with $N_{j-1}\ge1$. The conditional law of $y_{j-1}$ given $\alpha\in C$ has a density on $R_{j-1}(C)$ whose
supremum and infimum differ by a factor at most $(1+1/N_{j-1})^2\le1+3/N_{j-1}$.
\end{lemma}

\begin{proof}
With $t=\alpha_{j-1}=(N_j+y_{j-1})/N_{j-1}$ we have $|d\alpha/dy_{j-1}|=N_{j-1}^{-1}(q_{j-1}+q_{j-2}t)^{-2}$, and $t$ varies in an
interval of length $\le1/N_{j-1}$ on $C$; since $q_{j-2}\le q_{j-1}$ the ratio is at most $(1+1/N_{j-1})^2$.
\end{proof}

\begin{lemma}[Oscillations]\label{lem:osc}
\textup{(a)} Let $0\le k\le i-2$ and let $C$ be a level-$i$ cell with $i\le J$ on $C$. Then
\[
\osc_Cy_k\le e^4\sup_C(\alpha_k\alpha_{k+1}\cdots\alpha_{i-2})\le e^4\min\big(2^{-\lfloor (i-1-k)/2\rfloor},\,a_{k+1}^{-1}2^{-\lfloor(i-2-k)/2\rfloor}\big).
\]
If $i>J$ and $k\le J-2$, $\osc_Cy_k$ is bounded by the same expression with $i$ replaced by $J$.
\textup{(b)} Let $C$ be a level-$i$ cell, $i\ge1$, with $N_{i-1}\ge1$, on which $a_i,\dots,a_{l+1}$ are constant ($l\ge i-1$). Then
$\osc_C\alpha_l\le 4\,Q^2/N_{i-1}$, where $Q=\prod_{t=i}^{l}(a_t+1)$ (empty product $=1$).
\end{lemma}

\begin{proof}
(a) On $C$ all $N_1,\dots,N_i$ are fixed, $y_{i-1}$ ranges in an interval of length $\le1$, and by Lemma~\ref{lem:fibre} (applied at
levels $l=k+1,\dots,i-1$) $y_k$ is a smooth monotone function of $y_{i-1}$ with
$|dy_k/dy_{i-1}|=\prod_{l=k}^{i-2}\alpha_l(1+y_l/N_{l+1})$. Since $N_{l+2}\le N_l/2$, $\sum_{l>k}N_l^{-1}\le4/N_{i-1}\le4$.
The product of $r$ consecutive $\alpha$'s is at most $2^{-\lfloor r/2\rfloor}$ and $\alpha_k<1/a_{k+1}$. The case $i>J$ follows
because cells decrease with the level.
(b) On $C$, $\alpha_{i-1}$ varies in an interval of length $\le1/N_{i-1}$ (Lemma~\ref{lem:fibre}); writing
$\alpha_{i-1}=\frac{P+P'\alpha_l}{Q_0+Q_0'\alpha_l}$ with the convergent denominators $Q_0'\le Q_0\le Q$ of
$[0;a_i,\dots,a_l]$, we get $|d\alpha_l/d\alpha_{i-1}|=(Q_0+Q_0'\alpha_l)^2\le4Q^2$.
\end{proof}

\begin{remark}\label{rem:typechain}
By Lemma~\ref{lem:fibre}, for regular cells the children of a cell of type \textbf F, \textbf P, \textbf Q at level $j$ have types whose
conditional probabilities are, up to relative errors of order $(a_j+1)(a_{j+1}+1)/N_j$, those of a three-state Markov chain with
transition probabilities (with $x=\alpha_{j-1}$, $a=a_j$, $y=\alpha_j$)
$\textbf F\to\textbf F: ax$, $\textbf F\to\textbf P: xy$; $\textbf P\to\textbf P: y$, $\textbf P\to\textbf Q:1-y$;
$\textbf Q\to\textbf F:\frac{(a-1)x}{1-x}$, $\textbf Q\to\textbf P:\frac{xy}{1-x}$. Its three-step Dobrushin coefficient can be arbitrarily close
to $1$ (digits $(1,A,1,\dots)$ with $A$ large), although numerically the coefficients over $n\ge4$ steps appear to be bounded by $1/F_n$.
The proof below does not use this chain: it uses only the induced law of $y$, which is described by the scalar recursion of
Section~\ref{sec:transfer} with an explicit two-step contraction.
\end{remark}
\section{The two-piece transfer}\label{sec:transfer}

For $x\in(0,1)$ and $p\in[0,1]$ let $\nu_{x,p}$ be the probability measure on $[0,1)$ with density
$\frac{1-p}{1-x}$ on $[0,1-x)$ and $\frac px$ on $[1-x,1)$ (the \emph{$x$-two-piece law with upper mass $p$}); $\nu_{x,x}$ is
Lebesgue measure. For $x$ with $1/x\notin\Z$ write $a=\lfloor1/x\rfloor$, $y=\fr{1/x}$, $T_x(u)=\fr{-u/x}$, and
\begin{equation}\label{eq:lambdamu}
\lambda(x)=-\frac{y(1-y)}{a+y-1},\qquad \mu(x)=\frac{ay}{a+y-1}.
\end{equation}

\begin{lemma}[Exact transfer]\label{lem:exact}
$T_x{}_*\,\nu_{x,p}=\nu_{y,\ \lambda(x)p+\mu(x)}$.
\end{lemma}

\begin{proof}
Let $A=\frac{1-p}{1-x}$, $B=\frac px$. The density of the image is $g(v)=x\sum_cf(x(c-v))$, the sum over the integers $c$ with
$x(c-v)\in[0,1)$. Note $x(a-1)\le1-x\le xa$. For $v\in(0,1-y)$ the admissible $c$ are $1,\dots,a$; for $c\le a-1$,
$x(c-v)<1-x$, while for $c=a$, $x(a-v)>x(a-1+y)=1-x$. For $v\in(1-y,1)$ the admissible $c$ are $1,\dots,a+1$; now
$x(a-v)<1-x$ and $x(a+1-v)\in(xa,1)\subset[1-x,1)$. Hence $g=x((a-1)A+B)$ on $(0,1-y)$ and $g=x(aA+B)$ on $(1-y,1)$, so the image is
$y$-two-piece with upper mass $yx(aA+B)=y\big(\frac{a(1-p)}{a+y-1}+p\big)=\mu(x)+\lambda(x)p$, using $x=\frac1{a+y}$.
\end{proof}

Let $(x_l)_{l\ge l_0}$ be a Gauss orbit, $x_{l+1}=G(x_l)$, with digits $b_{l+1}=\lfloor1/x_l\rfloor$, and write
$\lambda_{l+1}=\lambda(x_l)$, $\mu_{l+1}=\mu(x_l)$ (these depend on $b_{l+1}$ and $x_{l+1}$). For $p_{l_0}\in[0,1]$ let
\begin{equation}\label{eq:rec}
p_{l+1}=\lambda_{l+1}p_l+\mu_{l+1}\qquad(l\ge l_0).
\end{equation}

\begin{lemma}[Uniform contraction]\label{lem:contr}
For all $l\ge l_0$:
\begin{enumerate}
\item[(i)] $p_{l+1}\in[0,1]$ and $1-p_{l+1}\le1-x_{l+1}$;
\item[(ii)] $p_{l}\le2x_{l}$ for $l\ge l_0+2$;
\item[(iii)] $|\lambda_{l+1}|\le1-x_{l+1}$, $|\lambda_{l+1}\lambda_{l+2}|\le3-2\sqrt2<\frac12$ and $|\lambda_{l+1}\lambda_{l+2}|\le x_{l+2}$;
\item[(iv)] if $(p'_l)$ is another solution of \eqref{eq:rec} with the same $(x_l)$, then for $m>l\ge l_0$
\[
|p_m-p'_m|\le 2^{-\lfloor(m-l)/2\rfloor}|p_l-p'_l|,\qquad\text{and}\qquad |p_m-p'_m|\le x_m\,2^{-\lfloor(m-l-2)/2\rfloor}|p_l-p_l'|\ \ (m\ge l+2).
\]
\end{enumerate}
\end{lemma}

\begin{proof}
Write $y=x_{l+1}$, $b=b_{l+1}$, $\kappa=\frac b{b+y-1}\ge1$, so $p_{l+1}=y\kappa(1-p_l)+yp_l$ and
$1-p_{l+1}=(1-p_l)(1-y\kappa)+p_l(1-y)$ with $1-y\kappa=\frac{(b-1)(1-y)}{b+y-1}\in[0,1-y]$; this gives (i).
For (ii): if $b_l\ge2$ then $\kappa\le2$ and $p_l\le x_l\max(\kappa,1)\le2x_l$. If $b_l=1$ then $p_l=(1-p_{l-1})+x_lp_{l-1}\le(1-x_{l-1})+x_l$
by (i), and $1-x_{l-1}=\frac{x_l}{1+x_l}\le x_l$.
For (iii): $|\lambda_{l+1}|=\frac{y(1-y)}{b+y-1}\le1-y$ because $b\ge1$. Write $z=x_{l+2}$. If $b_{l+2}\ge2$ then
$|\lambda_{l+2}|\le\frac{z(1-z)}{1+z}$ and $|\lambda_{l+1}|\le1$; if $b_{l+2}=1$ then $|\lambda_{l+1}|\le1-x_{l+1}=\frac{z}{1+z}$ and $|\lambda_{l+2}|\le1-z$.
In both cases $|\lambda_{l+1}\lambda_{l+2}|\le\frac{z(1-z)}{1+z}\le\min(z,\,3-2\sqrt2)$, the maximum of $\frac{z(1-z)}{1+z}$ on $(0,1)$ being attained
at $z=\sqrt2-1$.
(iv) follows from $p_m-p'_m=\big(\prod_{t=l+1}^m\lambda_t\big)(p_l-p'_l)$, grouping the factors in pairs (the last pair by the third
inequality of (iii)).
\end{proof}

For later use we record two elementary properties of $\nu_{x,p}$: its total variation distance to Lebesgue measure is $|p-x|$, and
its density is bounded when $p$ is comparable to $x$:
\begin{equation}\label{eq:nudens}
\|\nu_{x,p}-\nu_{x,x}\|_{TV}=|p-x|,\qquad \frac{d\nu_{x,p}}{du}\le\max\Big(\frac{1-p}{1-x},\frac px\Big)\le2\quad\text{if }p\le 2x\text{ and }1-p\le1-x .
\end{equation}
Moreover $\frac{\partial}{\partial p}\int g\,d\nu_{x,p}=\frac1x\int_{1-x}^1g-\frac1{1-x}\int_0^{1-x}g$, so $\big|\frac{\partial}{\partial p}\int g\,d\nu_{x,p}\big|\le2\sup|g|$.

\subsection*{The conditional law of the fibre variables}
For a level-$i$ cell $C$ which is $k$-regular ($i\le k$) fix a reference point $\alpha^*\in C$, put $x_l=\alpha_l(\alpha^*)$ and
$b_l=a_l$ (constant on $C$ for $l\le k+2$), and let $p^C_l$ ($i-1\le l\le k$) be the solution of \eqref{eq:rec} with
\[
p^C_{i-1}=x_{i-1},\ 1,\ 0\qquad\text{according as $C$ has type \textbf F, \textbf P, \textbf Q}.
\]
Define the error parameter
\begin{equation}\label{eq:EC}
E_C=\frac1{\min_CN_{k+1}}+\frac{1}{N_{i-1}}\prod_{u=i}^{k+2}(b_u+1)^5 .
\end{equation}

\begin{theorem}[Conditional law of $y_l$]\label{thm:condlaw}
There is an absolute constant $C_0$ such that, for every $k$-regular level-$i$ cell $C$ and every $i-1\le l\le k$,
\[
\big\|\mathrm{Law}(y_l\mid C)-\nu_{x_l,\,p^C_l}\big\|_{TV}\le C_0\,(l-i+2)\,E_C .
\]
\end{theorem}

\begin{proof}
Write $o_l=\osc_C\alpha_l$; by Lemma~\ref{lem:osc}(b), $o_l\le4\prod_{u=i}^{l}(b_u+1)^2/N_{i-1}$. Note $1/x_l\le b_{l+1}+1$ and
$1/(1-x_l)\le2(b_{l+2}+1)$. We may assume $C_0E_C\le1$, so that $o_l\le x_l/2$ and $o_l\le(1-x_l)/2$ for $l\le k$; the cell
$C$ has finitely many sub-cells, and we write $\underline N_l=\min_CN_l\ge\min_CN_{k+1}$ for $l\le k+1$.

\emph{Step 1.} Let $i-1\le l\le k$. The level-$(l+1)$ sub-cells $C'$ of $C$ partition $C$; by Lemma~\ref{lem:distortion} and
Proposition~\ref{prop:types}, $\mathrm{Law}(y_l\mid C')$ is within total variation $3/N_l$ of the uniform law on an interval of type
\textbf F, $(1-\alpha_l(e'),1)$ or $[0,1-\alpha_l(e'))$ with $|\alpha_l(e')-x_l|\le o_l$. Replacing $\alpha_l(e')$ by $x_l$ costs at most
$o_l/x_l+o_l/(1-x_l)$ in total variation. The resulting mixture is $x_l$-two-piece. Hence
$\|\mathrm{Law}(y_l\mid C)-\nu_{x_l,\tilde p_l}\|_{TV}\le e_l:=3/\underline N_l+o_l(x_l^{-1}+(1-x_l)^{-1})$ for some $\tilde p_l$, and
$\pi_l:=\Prob(y_l\ge1-x_l\mid C)$ satisfies $|\pi_l-\tilde p_l|\le e_l$. For $l=i-1$ the partition is trivial and the type of $C$ gives
$|\pi_{i-1}-p^C_{i-1}|\le 2e_{i-1}$.

\emph{Step 2.} Let $i-1\le l<k$. On $C$, $y_{l+1}=\fr{-y_l/\alpha_l}$ and $|y_l/\alpha_l-y_l/x_l|\le\eta_l:=2o_l/x_l^2$, so on the circle
$\R/\Z$ the points $y_{l+1}$ and $T_{x_l}(y_l)$ are within $\eta_l$. Thus $\{y_{l+1}\ge1-x_{l+1}\}$ and
$\{T_{x_l}(y_l)\ge1-x_{l+1}\}$ differ only if $T_{x_l}(y_l)$ is within $\eta_l$ of $0$ or of $1-x_{l+1}$. By Lemma~\ref{lem:exact}, under
$\nu_{x_l,\tilde p_l}$ the variable $T_{x_l}(u)$ has law $\nu_{x_{l+1},\lambda_{l+1}\tilde p_l+\mu_{l+1}}$, whose density is at most
$1/\min(x_{l+1},1-x_{l+1})$. Therefore
\[
\pi_{l+1}=\lambda_{l+1}\pi_l+\mu_{l+1}+r_l,\qquad |r_l|\le3e_l+\frac{4\eta_l}{\min(x_{l+1},1-x_{l+1})}.
\]

\emph{Step 3.} Since $|\lambda_t|\le1$, $|\pi_l-p^C_l|\le2e_{i-1}+\sum_{t=i-1}^{l-1}|r_t|$, and
$\|\mathrm{Law}(y_l\mid C)-\nu_{x_l,p^C_l}\|_{TV}\le e_l+|\tilde p_l-\pi_l|+|\pi_l-p^C_l|$ (note
$\|\nu_{x,p}-\nu_{x,p'}\|_{TV}=|p-p'|$). Inserting the bounds for $o_l$, $1/x_l$, $1/(1-x_l)$ and $\underline N_l\ge\min_CN_{k+1}$ gives the claim.
\end{proof}

\begin{corollary}[Forgetting]\label{cor:forget}
Let $C$ be a $k$-regular level-$i$ cell and $C'\subset C$ a level-$(i+1)$ sub-cell, $i+2\le k$. Using a common reference point
$\alpha^*\in C'$,
\[
|p^C_k-p^{C'}_k|\le x_k\,2^{-\lfloor(k-i-2)/2\rfloor}.
\]
Moreover $1-p^C_l\le1-x_l$ for $i\le l\le k$, and $p^C_l\le2x_l$ for $i+1\le l\le k$; in particular $\nu_{x_l,p^C_l}$ has density at most $2$ for
$i+1\le l\le k$.
\end{corollary}

\begin{proof}
$p^C$ and $p^{C'}$ solve \eqref{eq:rec} with the same environment from level $i$ on, with initial values in $[0,1]$ at level $i$, so the first claim is
Lemma~\ref{lem:contr}(iv). The others are Lemma~\ref{lem:contr}(i)--(ii) and \eqref{eq:nudens}.
\end{proof}

\begin{remark}
Theorem~\ref{thm:condlaw} expresses the conditional law of $y_k$ given a cell at level $i$ through the \emph{digits only}, up to the
initial value $p_{i-1}^C$, whose influence decays like $2^{-(k-i)/2}$ \emph{uniformly in the digits}. Section~\ref{sec:cells} and the
present section are taken from \cite{E1002}, where the recursion is also checked numerically.
\end{remark}
\section{Facts on continued fractions}\label{sec:cf}

We use the following standard facts (see \cite{Kh64,IK02}).

\smallskip\noindent\textbf{(CF1) Cylinders.} For integers $b_1,\dots,b_m\ge1$ with denominators $q_t$,
$\lambda\{a_1=b_1,\dots,a_m=b_m\}=\frac{1}{q_m(q_m+q_{m-1})}\le\prod_{t=1}^mb_t^{-2}$. The density $h_j$ of $G^j_*\lambda$ satisfies
$\frac12\le h_j\le2$: indeed $h_j(x)=\sum(q_j+q_{j-1}x)^{-2}$, the sum running over cylinders of depth $j$, and
$\frac{1}{2q_j(q_j+q_{j-1})}\le(q_j+q_{j-1}x)^{-2}\le\frac{2}{q_j(q_j+q_{j-1})}$ because $q_{j-1}\le q_j$.

\smallskip\noindent\textbf{(CF2) Mixing (Kuzmin--L\'evy--Wirsing, $\psi$-mixing).} This is the Gauss--Kuzmin--Wirsing theorem applied to the
normalised restrictions of $\lambda$ to cylinders, whose images under $G^j$ have densities $\frac{1+u}{(1+ut)^2}$, $u=q_{j-1}/q_j\in[0,1]$, a
family on which the convergence to the Gauss density is uniform and exponential \cite[Ch.~2]{IK02}. There are constants $C_1$ and $\rho\in(0,1)$
such that for all $j,n\ge0$, all $A\in\sigma(a_1,\dots,a_j)$ and all Borel $E\subset(0,1)$,
\[
\big|\lambda(A\cap G^{-(j+n)}E)-\lambda(A)\gamma(E)\big|\le C_1\rho^n\lambda(A)\gamma(E).
\]
Since the densities $h_j$ of $G^j_*\lambda$ lie in $[\frac12,2]$ (by the computation in (CF1)), it follows that there is a constant
$C_2$ such that, if $f$ is $\sigma(a_1,\dots,a_j)$-measurable and $g=\hat g\circ G^{j+n}$ for a Borel function $\hat g$, then
\[
|\E[fg]-\E f\,\E g|\le C_2\rho^n\,\E|f|\,\E|g|,\qquad \Big|\E g-\int\hat g\,d\gamma\Big|\le C_1\rho^{j+n}\int|\hat g|\,d\gamma ;
\]
in particular $\Prob(a_{k+1}=b)=\frac{1+O(\rho^k)+O(1/b)}{b^2\log2}$.

\smallskip\noindent\textbf{(CF3) Products of digits.} For $0<\delta\le\frac12$, $j\ge0$ and $m\ge1$,
$\E\prod_{t=j+1}^{j+m}(a_t+1)^{1-\delta}\le2(4/\delta)^m$; hence $\Prob\big(\prod_{t=j+1}^{j+m}(a_t+1)>x\big)\le2(4/\delta)^mx^{-1+\delta}$.

\begin{proof}
By (CF1) applied to $G^j_*\lambda$, the expectation is at most
\[2\Big(\sum_{b\ge1}(b+1)^{1-\delta}b^{-2}\Big)^m\le2\Big(2\sum_bb^{-1-\delta}\Big)^m\le2(4/\delta)^m.\]
\end{proof}

\smallskip\noindent\textbf{(CF4) Growth of $q_k$.} $\E\log q_k=\frac{\pi^2}{12\log2}k+O(1)$ and $\mathrm{Var}(\log q_k)\le C k$ (L\'evy; for the
variance bound write $\log q_k=\sum_{j<k}\log(1/\alpha_j)+O(1)$ by \eqref{eq:theta} and replace $\log(1/\alpha_j)$ by its conditional
expectation given $a_{j+1},\dots,a_{j+r}$, which changes it by $O(2^{-r})$; (CF2) then bounds the covariances. In fact $\log q_k$ satisfies a
central limit theorem \cite{IK02}). We write
\[
c_0=\frac{12\log2}{\pi^2}.
\]

\smallskip\noindent\textbf{(CF5) Large digits.} $\Prob(a_{k+1}\ge b)\le 2/b$ for all $k\ge0$, $b\ge1$, since $h_k\le2$.

\smallskip\noindent\textbf{(CF6) The next digit.} For $k\ge0$ and $b\ge1$,
\[
\Prob(a_{k+1}=b\mid\cF_k)=p_{\omega_k}(b):=\frac{1+\omega_k}{(b+\omega_k)(b+1+\omega_k)},\qquad \Prob(a_{k+1}>T\mid\cF_k)=\frac{1+\omega_k}{\lfloor T\rfloor+1+\omega_k}\le\frac2T ,
\]
where $\omega_k=q_{k-1}/q_k$, because conditionally on $\cF_k$ the variable $\alpha_k$ has density $(1+\omega_k)(1+\omega_kt)^{-2}$ on $(0,1)$ (the case
$\cI=(0,1)$ of Lemma~\ref{lem:refcyl}). In particular $|p_\omega(b)-(1+\omega)b^{-2}|\le10b^{-3}$ and $|\partial_\omega p_\omega(b)|\le 3b^{-2}$.

\smallskip
Functions of the digits alone are handled by (CF2); functions which also involve the residue vectors $w_k$ or the endpoint orbits are handled by
the mixing of the extended chain, Lemmas~\ref{lem:mixing} and~\ref{lem:statlaw}.
\section{Set-up and decomposition}\label{sec:decomp}

\subsection{Parameters}
Let $N\to\infty$, $L=\log N$, and fix $\eps\in(0,\frac12)$. Put
\[
T=\eps L,\quad U=L^{1.1},\quad Z=L^{10},\quad s=\lceil\log\log L\rceil,\quad M=\lceil 3L\rceil,\quad K_0=\lfloor c_0(L-10\log L)\rfloor .
\]
For large $N$ we have $J\le M$ (Lemma~\ref{lem:sizes} and $q_j\ge 2^{(j-1)/2}$). We write $o_\Prob(1)$ for random variables tending to
$0$ in probability as $N\to\infty$ ($\eps$ fixed), and $\Pi_k=\prod_{u=k-s}^{k+2}(a_u+1)$.

\subsection{Bulk indicators}
For $s<k\le M$ define the $\cH_{k-s}$-measurable indicator
\[
g_k=\ind\{N_{k-s}\ge Z\}\cdot\ind\{\cC_{k-s}(\alpha)\text{ is $k$-regular}\}\cdot\ind\{\Pi_k\le L^{1.2}\}
\]
(on a $k$-regular level-$(k-s)$ cell the digits $a_{k-s},\dots,a_{k+2}$ are constant, so $g_k$ is a function of the cell), and split
$g_k=\chi_k+\chi_k'$ with $\chi_k=g_k\ind\{a_{k+1}\le T\}$, $\chi'_k=g_k\ind\{a_{k+1}>T\}$; both are $\cH_{k-s}$-measurable.

\begin{lemma}\label{lem:bulkfacts}
If $g_k=1$ then $N_l\ge L^{8.7}$ for all $l\le k+1$, in particular $k+1<J$, and every level-$i$ cell $C\subset\cC_{k-s}(\alpha)$ with
$k-s\le i\le k$ is $k$-regular with $E_C\le2L^{-2.7}$ (see \eqref{eq:EC}).
\end{lemma}

\begin{proof}
$N_{l+1}\ge N_l\alpha_l-1\ge N_l/(a_{l+1}+1)-1$, so $N_{k+1}\ge Z/\Pi_k-2\ge L^{8.7}$ and $N_l\ge N_{k+1}$ for $l\le k+1$. Sub-cells of a
$k$-regular cell are $k$-regular (Section~\ref{sec:cells}), and $E_C\le N_{k+1}^{-1}+\Pi_k^5N_{i-1}^{-1}\le 2L^{6-8.7}$.
\end{proof}

\begin{lemma}[Good event]\label{lem:good}
Let $G_N$ be the event that $a_t\le U$ for all $t\le M+3$, that $\Pi_k\le L^{1.2}$ for all $s<k\le M$, and that $\cC_{k-s}(\alpha)$ is
$k$-regular for every $s<k\le M$ with $N_{k-s}\ge Z$. Then $\Prob(G_N)\to1$.
\end{lemma}

\begin{proof}
By (CF5), $\Prob(\exists t\le M+3:a_t>U)\le2(M+3)/U\to0$. By (CF3) with $\delta=\frac1{10}$ and $m=s+3$,
$\Prob(\Pi_k>L^{1.2})\le2\cdot40^{s+3}L^{-1.08}=L^{-1.08+o(1)}$, and we sum over $k\le M$.
For regularity, fix $k$ with $N_{k-s}\ge Z$ and $\Pi_k\le L^{1.2}$, let $C=\cC_{k-s}(\alpha)$ and let $D=D_{k+2}(\alpha)$ be the cylinder of depth
$k+2$ containing $\alpha$. \emph{Case 1:} $\overline C\subset\mathrm{int}\,D$. Then $a_1,\dots,a_{k+2}$ are constant on a neighbourhood of
$\overline C$; in particular $\Pi_k$ is constant on $C$, so the bound $N_{l+1}\ge N_l/(a_{l+1}+1)-1$ of Lemma~\ref{lem:bulkfacts} gives
$N_l\ge L^{8.7}\ge2$ on all of $C$ for $l\le k+1$, and $C$ is $k$-regular. \emph{Case 2:} $\overline C\not\subset\mathrm{int}\,D$. Then
$\mathrm{dist}(\alpha,\partial D)\le|C|$. By the proof of Lemma~\ref{lem:distortion} and Lemma~\ref{lem:sizes},
$|C|\le(N_{k-s-1}q_{k-s-1}^2)^{-1}\le4/(Nq_{k-s-1})$, while $\lambda(D)\ge(2q_{k+2}^2)^{-1}$ and $q_{k+2}\le q_{k-s-1}\Pi_k$. Since $N_{k-s}\ge Z$
forces $q_{k-s-1}<q_{k-s}<N/Z$, we get $|C|/\lambda(D)\le8\Pi_k^2/Z\le L^{-3}$. So Case~2 implies
$\mathrm{dist}(\alpha,\partial D_{k+2}(\alpha))\le L^{-3}\lambda(D_{k+2}(\alpha))$. Since the cylinders of depth $k+2$ partition $(0,1)$ up to a
null set, $\lambda\{\alpha:\mathrm{dist}(\alpha,\partial D_{k+2}(\alpha))\le\eta\lambda(D_{k+2}(\alpha))\}=2\eta$ for $\eta\le\frac12$; with $\eta=L^{-3}$ and a
union bound over $k\le M$, the probability of non-regularity is $O(L^{-2})$.
\end{proof}

\begin{lemma}[Bulk end]\label{lem:bulkend}
Let $K''=\max\{l:N_l\ge Z\}$. Then $J-K''\le2\log_2(Z+4)+4$, and for every $\eta>0$ there is $A$ such that
$\Prob(|K''-K_0|>A\sqrt L)\le\eta$ for all large $N$.
\end{lemma}

\begin{proof}
By Lemma~\ref{lem:sizes}, $N|\theta_{K''}|<Z+4$; for $K''<k<J$ we need $N_k\ge1$, i.e.\ $N|\theta_{k-1}|\ge1$, and $|\theta_{k+1}|\le|\theta_{k-1}|/2$, which gives the first
claim. For the second, $N_l\ge Z$ holds if $q_l\le N/(2Z+8)$ and fails if $q_l>N/Z$ (Lemma~\ref{lem:sizes}). Hence $m_1\le K''\le m_2$ with
$m_1=\max\{l:q_l\le N/(2Z+8)\}$ and $m_2=\max\{l:q_l\le N/Z\}$. As $q_l$ increases with $l$, $\{m_2>K_0+A\sqrt L\}\subset\{\log q_{l_+}\le L-10\log L\}$
with $l_+=K_0+\lceil A\sqrt L\rceil$, and $\{m_1<K_0-A\sqrt L\}\subset\{\log q_{l_-}>L-10\log L-\log(2+8/Z)\}$ with $l_-=K_0-\lfloor A\sqrt L\rfloor$.
By (CF4), $\E\log q_{l_\pm}=L-10\log L\pm A\sqrt L/c_0+O(1)$ and $\mathrm{Var}\log q_{l_\pm}=O(L)$, so Chebyshev's inequality bounds both
probabilities by $Cc_0^2/A^2$, which is $\le\eta/2$ for $A$ large.
\end{proof}

\subsection{The decomposition}
From now on $\beta,r\in\Q$ are fixed as in Theorem~\ref{thm:main}, with endpoints $z_1=\fr{-\beta}$, $z_2=\fr{r-\beta}$ and common denominator $Q$;
constants may depend on $Q$ (hence on $\beta,r$). If one of the endpoints is $0$, all terms attached to it below are $0$. For an endpoint $z_i\ne0$
write $z_{i,k}$ for its orbit, $u_{i,k}=\fr{P_i(w_k)_2/Q}$ with $P_i=Qz_i$, and $v_{i,k}$ as in Lemma~\ref{lem:orbit}; for $z_i=0$ put $u_{i,k}=0$.
Writing $z'=\fr{-z/x}$ and $\tilde\kappa=\kappa-z'$, i.e.
\[
\tilde\kappa(x,y,z)=\begin{cases}-1 & y\ge z,\\ -(1-y')+\ind\{y'<z'\}-z' & y<z,\end{cases}
\]
we have $\ell(x,y,z)=\tau(y,z)/x+z'+\tilde\kappa(x,y,z)$. Since $y'-z'+\ind\{y'<z'\}=\fr{y'-z'}$ for $y',z'\in[0,1)$,
\begin{equation}\label{eq:kappaid}
\tilde\kappa(x,y,z)=-1+\ind\{y<z\}\,\fr{y'-z'}\in[-1,0).
\end{equation}
Hence $\tilde\kappa$ depends on $(y',z')$ only through $y'-z'$ modulo $1$; it is continuous in $z'$ on the circle and in $y$ at $y'=0$, it is continuous at
$y=z$ (as $y\uparrow z$, $y'\downarrow z'$ and $\fr{y'-z'}\to0$), and its only other discontinuities are jumps of size $1$ where $y'=z'$ with $y<z$, or where
the case $y<z$ changes because $z$ passes through $0$. Accordingly, by \eqref{eq:twoexp},
\begin{equation}\label{eq:LMZK}
\begin{gathered}
L_k=M_k+Z_k+K_k,\qquad M_k=\frac{\tau(y_k,z_{2,k})-\tau(y_k,z_{1,k})}{\alpha_k},\qquad Z_k=z_{2,k+1}-z_{1,k+1},\\
K_k=\tilde\kappa(\alpha_k,y_k,z_{2,k})-\tilde\kappa(\alpha_k,y_k,z_{1,k}),
\end{gathered}
\end{equation}
with $|M_k|\le\frac{a_{k+1}+1}4$, $|Z_k|\le1$, $|K_k|\le1$. The variable $Z_k$ depends on $\alpha$ only. We split
\begin{equation}\label{eq:decomp}
D_N(\alpha,\beta;r)=\Lambda+\Lambda_Z+H+R,
\end{equation}
\[
\Lambda=\sum_{k=s+1}^{M}(-1)^k\chi_k(M_k+K_k),\quad \Lambda_Z=\sum_{k=s+1}^{M}(-1)^k\chi_kZ_k,\quad H=\sum_{k=s+1}^{M}(-1)^k\chi'_kL_k,
\]
and $R=\sum_{k<J}(-1)^kL_k(1-g_k)$, with $g_k:=0$ for $k\le s$.

\begin{lemma}\label{lem:R}
$R/L\to0$ in probability.
\end{lemma}

\begin{proof}
Recall $|L_k|\le2(a_{k+1}+1)$. The terms with $k\le s$ contribute at most $2\sum_{k\le s}(a_{k+1}+1)$, which is $O_\Prob(s^2)$ by (CF5).
On $G_N$, $g_k=\ind\{N_{k-s}\ge Z\}$ for $s<k\le M$, so the remaining terms have $K''+s<k<J$; by Lemma~\ref{lem:bulkend} there are at most
$2\log_2(Z+4)+4=O(\log L)$ of them, and with probability $\ge1-\eta$ they lie in the deterministic window
$W=[K_0-A\sqrt L,\,K_0+A\sqrt L+C\log L]$. Hence, for $\delta>0$,
\[
\Prob(|R|>\delta L)\le\Prob(G_N^c)+\eta+\Prob\Big(\sum_{k\le s}(a_{k+1}+1)>\tfrac{\delta L}4\Big)+\Prob\Big(\max_{k\in W}a_{k+1}>\tfrac{\delta L}{C'\log L}\Big),
\]
and by (CF5) the last two terms are $O(s^2/(\delta L))$ and $O(\sqrt L\log L/(\delta L))$.
\end{proof}

Thus it remains to show that $(\Lambda+\Lambda_Z)/L$ is small in probability uniformly as $\eps\to0$ (Section~\ref{sec:light}) and that $H/L$
converges to a compound Poisson law (Section~\ref{sec:heavy}).
\section{The bounded part}\label{sec:light}

\begin{proposition}\label{prop:light}
There is a constant $C_Q$ such that for every $\delta>0$, $\limsup_{N\to\infty}\Prob(|\Lambda+\Lambda_Z|>\delta L)\le C_Q\eps/\delta^2$.
\end{proposition}

Throughout this section we work on $G_N\cap\{|K''-K_0|\le A\sqrt L\}$, whose complement has probability $\le\eta+o(1)$ with $\eta>0$ arbitrary
(Lemmas~\ref{lem:good} and~\ref{lem:bulkend}); there $\chi_k=\ind\{a_{k+1}\le T\}\ind\{k\le K''+s\}$ for $s<k\le M$. Let
$W'=[K_0-A\sqrt L,\,K_0+A\sqrt L+s]$, a deterministic window of the same kind as in Lemma~\ref{lem:R}.

\subsection{Approximation by finitely many pieces}
For $l\ge0$ and $r\ge1$ the pieces $\fp_l,\dots,\fp_{l+r}$ determine the digits $a_{l+1},\dots,a_{l+r+1}$, so $\alpha_l$ is determined up to an interval of
length at most $F_{r+1}^{-2}$ ($F_j$ the Fibonacci numbers). We write $\alpha_l^{[r]}$ for the midpoint of this interval; it is a function of
$\fp_l,\dots,\fp_{l+r}$ and $|\alpha_l-\alpha_l^{[r]}|\le F_{r+1}^{-2}$. Inside a fixed piece $\fp_l$ the integer parts $n^{(i)}_l$ are constant
(Lemma~\ref{lem:pieces}(a)), so $z_{i,l}=u_{i,l}+v_{i,l}\alpha_l-n^{(i)}_l$ is an affine function of $\alpha_l$ with slope $|v_{i,l}|<4$.

\subsection{The part depending on $\alpha$ only}
On the event above, $\Lambda_Z=\Sigma_Z+O(|W'|)$ with $\Sigma_Z=\sum_{s<k\le K_0}(-1)^k\hat z_k$ and $\hat z_k=\ind\{a_{k+1}\le T\}Z_k$; note $|W'|=o(L)$.

\begin{lemma}\label{lem:sigmaZ}
$\E[\Sigma_Z^2]\le C_QL$.
\end{lemma}

\begin{proof}
$\hat z_k=\ind\{a(\fp_k)\le T\}\big(z_{2,k+1}-z_{1,k+1}\big)$ is a functional of the chain from time $k+1$ together with $\fp_k$, i.e.\ of the chain from time $k$,
the same functional for every $k$, bounded by $1$ and independent of $\omega$. By Lemma~\ref{lem:mixing}(b), $\E\hat z_k=\tilde\mu+O(\vartheta_Q'^{\,k})$, so the
alternating sum of the means is $O(1)$. Let $\hat z^{[r]}_k$ be obtained by replacing $\alpha_{k+1}$ by $\alpha^{[r]}_{k+1}$; then
$|\hat z_k-\hat z_k^{[r]}|\le8F_{r+1}^{-2}$ and $\hat z^{[r]}_k$ is $\cF'_{k+r+2}$-measurable. For $j<k$ and $r=\lfloor(k-j)/2\rfloor$, Lemma~\ref{lem:mixing}(a)
gives $|\mathrm{Cov}(\hat z_j,\hat z_k)|\le|\mathrm{Cov}(\hat z^{[r]}_j,\hat z_k)|+16F_{r+1}^{-2}\le C_Q\vartheta_Q'^{\,k-j-r-2}+16F_{r+1}^{-2}$, which is summable in
$k-j$. Hence $\mathrm{Var}(\Sigma_Z)\le C_QL$.
\end{proof}

\subsection{Continuity properties}
Call an interval $I$ of parameters \emph{$k$-clean} if $a_1,\dots,a_{k+2}$ are constant on $I$ and $\alpha_{k-1}(I)$, $\alpha_k(I)$ contain no point of
$E_0$. We use that $G$ maps $E_0\setminus\{1/a:a\ge1\}$ into $E_0$ (proof of Lemma~\ref{lem:pieces}(b)). Hence, on an interval on which $a_1,\dots,a_{k+2}$
are constant, a parameter at which $\alpha_j$ ($j\le k$) passes through a point $e\in E_0$ is also a parameter at which $\alpha_{j+1},\dots,\alpha_k$ pass
through the points $G(e),\dots,G^{k-j}(e)\in E_0$: these are not branch points, since $\alpha_{j+1},\dots,\alpha_k$ stay inside fixed open branches.

\begin{lemma}\label{lem:zcont}
On a $k$-clean interval $I$ the pieces $\fp_0,\dots,\fp_k$ are constant, hence so are $\zeta_k$ and $\zeta_{k+1}$. Consequently $z_{i,k}=u_{i,k}+v_{i,k}\alpha_k-n_{i,k}$ is a
continuous function on $I$ with values in $(0,1)$, $z_{i,k+1}=\fr{u_{i,k+1}+v_{i,k+1}\alpha_{k+1}}$ is continuous as a map to the circle $\R/\Z$, and
$|dz_{i,l}/d\alpha|\le16q_l^2$ for $l=k,k+1$.
\end{lemma}

\begin{proof}
By the remark above, a change of $\fp_j$ for some $j\le k$ on $I$ (which can only happen where $\alpha_j$ passes through a point of $E_0$, the digits
being constant) would force $\alpha_k$ to pass through a point of $E_0$, which is excluded. So $\fp_0,\dots,\fp_k$ are constant and, by \eqref{eq:zetarec},
so are $\zeta_k,\zeta_{k+1}$. The formulas follow from Lemma~\ref{lem:orbit}; $z_{i,k}$ does not vanish on $I$ because $\alpha_k(I)$ contains no point of $E_0$.
Finally $|dz_{i,l}/d\alpha|\le|v_{i,l}|\,|d\alpha_l/d\alpha|\le4(q_l+q_{l-1}\alpha_l)^2\le16q_l^2$.
\end{proof}

If $a_1,\dots,a_{k+2}$ are constant on an interval but it is not $k$-clean, consider a parameter at which some $\alpha_j$ passes through $E_0$. For each
endpoint, at most one of the circle points $z_{i,l}$ jumps there: if $z_{i,j}$ passes through $0$, then $z_{i,j+1}$ has the one-sided limits $1^-$ and
$(1-\alpha_{j+1})^+$, both are mapped by $t\mapsto\fr{-t/\alpha_{j+1}}$ to $1-\alpha_{j+2}$, so $z_{i,j+2}$ is continuous with value $1-\alpha_{j+2}\ne0$, and
inductively $z_{i,l}=1-\alpha_l\ne0$ is continuous for $l\ge j+2$. In particular $z_{i,k}$ jumps by at most $\alpha_k$ and $z_{i,k+1}$ by at most $\alpha_{k+1}$ on
the circle.

Let $G_k=M_k+K_k$. On a $k$-clean interval $G_k$ is a continuous function of $(\alpha_k,y_k,y_{k+1},z_{i,k},z_{i,k+1})$, except at parameters where
$y_k<z_{i,k}$ and $y_{k+1}=z_{i,k+1}$ for some $i$, where $K_k$ jumps by $1$: indeed $\tau$ is $1$-Lipschitz on the circle in each variable, and
$\tilde\kappa$ is continuous at $y=z$, at $y'=0$ and in $z'\in\R/\Z$, and $1$-Lipschitz in $(y',z')$ elsewhere \eqref{eq:kappaid}. Without the
cleanness assumption, every discontinuity of $G_k$ on an interval on which $a_1,\dots,a_{k+2}$ are constant is a jump of size at most $10$: at a parameter
where $\alpha_k$ passes through $E_0$, each $z_{i,k}$ moves by at most $\alpha_k$ on the circle, hence $M_k$ by at most $2$, and each $\tilde\kappa\in[-1,0)$
changes by at most $1$.

\subsection{Oscillation bounds}
For a level-$i$ cell $C$ let $\cB_k(C)$ be the event that $C$ is not $k$-clean or that $y_{k+1}-z_{i',k+1}$ vanishes (on the circle) at some point of $C$ for
some endpoint $i'$; it is a property of the cell.

\begin{lemma}\label{lem:oscG}
Suppose $\chi_k=1$, $k+3\le i\le J$, and let $C=\cC_i(\alpha)$. Then
\begin{equation}\label{eq:oscX}
\osc_CG_k\le 6e^4\,2^{-\lfloor(i-k-2)/2\rfloor}+L^{-5}+90\,\ind_{\cB_k(C)} ,
\end{equation}
and
\begin{equation}\label{eq:PB}
\Prob\big(\chi_k=1,\ i\le J,\ \cB_k(\cC_i(\alpha))\big)\le C\,2^{-\lfloor(i-k-2)/2\rfloor}+C_Q(s+2)L^{-2.7}.
\end{equation}
The same bounds hold with $i=J$ and $2^{-\lfloor(i-k-2)/2\rfloor}$ replaced by $L^{-8}$.
\end{lemma}

\begin{proof}
On $C$, $a_1,\dots,a_{k+2}$ and $N_1,\dots,N_i$ are constant ($C\subset\cC_{k+1}(\alpha)$, which is $k$-regular). By Lemma~\ref{lem:osc}(a),
$\osc_Cy_k\le e^4a_{k+1}^{-1}2^{-\lfloor(i-k-2)/2\rfloor}$ and $\osc_Cy_{k+1}\le e^42^{-\lfloor(i-k-2)/2\rfloor}$; by Lemma~\ref{lem:osc}(b), $\osc_C\alpha_k\le4/N_k$ and
$\osc_C\alpha_{k-1}\le4/N_{k-1}$. A level-$i$ cell has length $|C|\le(N_{i-1}q_{i-1}^2)^{-1}$ (proof of Lemma~\ref{lem:distortion}), and
$N_{i-1}q_{i-1}\ge N/16$ by Lemma~\ref{lem:sizes}; hence on a $k$-clean cell, by Lemma~\ref{lem:zcont}, $\osc_Cz_{i',l}\le256q_{k+1}/N\le256/N_{k+1}$ for
$l\le k+1$. Since $|\partial_y\tau|,|\partial_z\tau|\le1$, $\alpha_k^{-1}\le a_{k+1}+1\le T+1$, $\tau\le\frac14$, and $N_l\ge L^{8.7}$ for $l\le k+1$ (Lemma~\ref{lem:bulkfacts}),
\[
\osc_CM_k\le2(a_{k+1}+1)\osc_Cy_k+(T+1)\sum_{i'}\osc_Cz_{i',k}+(T+1)^2\osc_C\alpha_k\le4e^4\,2^{-\lfloor(i-k-2)/2\rfloor}+\tfrac12L^{-5},
\]
and, if moreover $y_{k+1}-z_{i',k+1}$ has no zero on $C$, $\osc_CK_k\le2\osc_Cy_{k+1}+2\sum_{i'}\osc_Cz_{i',k+1}\le2e^4\,2^{-\lfloor(i-k-2)/2\rfloor}+\tfrac12L^{-5}$.
Otherwise $G_k$ has at most $9$ discontinuities on $C$. Indeed, $\alpha_k(C)$ has length $<D_0^{-2}$, so it contains at most one point of $E_0$ (distinct
points of $E_0$ are at distance $\ge D_0^{-2}$), and by the remark before Lemma~\ref{lem:zcont} every parameter at which some $\alpha_j$, $j\le k$, passes
through $E_0$ is one at which $\alpha_k$ does; so there is at most one such parameter. On each of the (at most two) remaining subintervals $\zeta_{k+1}$ is
constant, and a lift of $y_{k+1}-z_{i',k+1}$ is strictly increasing as a function of $\alpha_{k+1}$ ($y_{k+1}=N_{k+1}\alpha_{k+1}-N_{k+2}$ increases, and
$z_{i',k+1}$ is affine in $\alpha_{k+1}$ with slope $v_{i',k+1}\le0$), with total increase $<1+256/N_{k+1}<2$; so it passes through an integer at most twice.
This gives at most $1+2\cdot2\cdot2$ discontinuities, each a jump of size at most $10$, which gives \eqref{eq:oscX}.

For \eqref{eq:PB}: if $C$ is not $k$-clean, then $\alpha_{k-1}$ or $\alpha_k$ is within $4L^{-8.7}$ of $E_0$ at the point $\alpha$; since the density of
$\alpha_l$ is at most $2$ (CF1), this has probability $\le C_QL^{-8.7}$. Otherwise a zero of $y_{k+1}-z_{i',k+1}$ on $C$ forces
$\|y_{k+1}-z_{i',k+1}\|\le\delta:=e^42^{-\lfloor(i-k-3)/2\rfloor}+256L^{-8.7}$ at $\alpha$. Condition on $C'=\cC_{k-s}(\alpha)$, which is $k$-regular with
$E_{C'}\le2L^{-2.7}$ (Lemma~\ref{lem:bulkfacts}), and fix a reference point $\alpha^*\in C'$. By Theorem~\ref{thm:condlaw}, $\mathrm{Law}(y_k\mid C')$ is within total
variation $C_0(s+2)E_{C'}$ of $\nu_{x_k,p}$, $p=p^{C'}_k$. On $C'$, $\osc\alpha_k\le4\Pi_k^2/N_{k-s-1}\le4L^{-7.6}$ (Lemma~\ref{lem:osc}(b)), so
$\|y_{k+1}-T_{x_k}(y_k)\|\le2\cdot4L^{-7.6}(T+1)^2\le L^{-5}$; by Lemma~\ref{lem:exact}, $T_{x_k}$ maps $\nu_{x_k,p}$ to $\nu_{x_{k+1},p'}$ with $p'=\lambda(x_k)p+\mu(x_k)$,
whose density is at most $2$ (Corollary~\ref{cor:forget}, the recursion being extended to level $k+1$ with the digit $a_{k+2}$, constant on $C'$). Finally,
off an event of probability $\le C_QL^{-7.6}$ (on which $C'$ is not $k$-clean), $z_{i',k+1}$ varies on $C'$ by at most
$16q_{k+1}^2|C'|\le256\Pi_k/N_{k+1}\le L^{-7}$. Hence $\Prob(\|y_{k+1}-z_{i',k+1}\|\le\delta\mid C')\le4(\delta+2L^{-5})+C_0(s+2)E_{C'}$, and \eqref{eq:PB}
follows. For $i=J$ use that $\chi_k=1$ forces $N_{k+1}\ge L^{8.7}$, hence $J-k\ge17\log_2L-2$ (as $N_{l+2}<N_l/2$), so all the factors
$2^{-\lfloor(i-k-2)/2\rfloor}$ above are at most $L^{-8}$.
\end{proof}

\subsection{The martingale decomposition}
Put $X_k=(-1)^k\chi_kG_k$ for $s<k\le M$, so $\Lambda=\sum_kX_k$, and $\Delta_{k,i}=\E[X_k\mid\cH_{i+1}]-\E[X_k\mid\cH_i]$. For the
$(\cH_i)$-stopping time $J\le M$ put $\E[X_k\mid\cH_J]=\sum_j\ind\{J=j\}\E[X_k\mid\cH_j]$. Since $\chi_k$ is $\cH_{k-s}$-measurable and
$\chi_k=1$ forces $J>k+1$, on $\{J\le k-s\}\in\cH_{k-s}$ all the conditional expectations below vanish, while on $\{J>k-s\}$ the
telescoping of the martingale $i\mapsto\E[X_k\mid\cH_i]$ between $k-s$ and $J$ gives, pointwise,
\[
X_k=\E[X_k\mid\cH_{k-s}]+\sum_{i=k-s}^{M-1}\ind\{i<J\}\Delta_{k,i}+\big(X_k-\E[X_k\mid\cH_J]\big).
\]
Summing over $k$ we obtain $\Lambda=\Lambda_{\rm res}+\Lambda_{\rm mart}+\Lambda_{\rm tail}$ with
\[
\Lambda_{\rm res}=\sum_k\E[X_k\mid\cH_{k-s}],\qquad \Lambda_{\rm mart}=\sum_{i<M}\ind\{i<J\}D_i,\quad D_i=\sum_{s<k\le\min(M,i+s)}\Delta_{k,i}.
\]

For a $k$-regular level-$i$ cell $C$ ($k-s\le i\le k$) with reference point $\alpha^*\in C$, let
\begin{gather*}
\bar G(x,p,z_1,z_2)=\int_0^1G(x,u,z_1,z_2)\,\nu_{x,p}(du),\\ G(x,u,z_1,z_2)=\frac{\tau(u,z_2)-\tau(u,z_1)}{x}+\tilde\kappa(x,u,z_2)-\tilde\kappa(x,u,z_1),
\end{gather*}
where inside $\tilde\kappa$ we use $u'=\fr{-u/x}$ and $z_i'=\fr{-z_i/x}$. Then $G_k=G(\alpha_k,y_k,z_{1,k},z_{2,k})$ and
$\sup_u|G(x,u,\cdot)|\le\frac{a+1}4+4$ with $a=\lfloor1/x\rfloor$; by the remark after \eqref{eq:nudens},
\begin{equation}\label{eq:Gbarp}
\Big|\frac{\partial\bar G}{\partial p}\Big|\le2(a+5).
\end{equation}

\begin{lemma}\label{lem:Gbarapprox}
If $\chi_k=1$ and $C=\cC_i(\alpha)$ with $k-s\le i\le k$, then, with the reference point $\alpha^*=\alpha$,
\[
\big|\E[G_k\mid C]-\bar G\big(\alpha_k,p^C_k,z_{1,k},z_{2,k}\big)\big|\le C(T+5)(s+2)L^{-2.7}+20\cdot\ind\{C\text{ is not $k$-clean}\},
\]
and $\Prob(\chi_k=1,\ \cC_{k-s}(\alpha)\text{ is not $k$-clean})\le C_QL^{-7.6}$.
\end{lemma}

\begin{proof}
By Theorem~\ref{thm:condlaw}, replacing $\mathrm{Law}(y_k\mid C)$ by $\nu_{x_k,p^C_k}$ costs at most $2(\frac{T+1}4+4)C_0(s+2)E_C$. On $C$ the parameters
$\alpha_k$ and (if $C$ is $k$-clean) $z_{i',k}$, $z_{i',k+1}$ vary by at most $4L^{-7.6}$ and $L^{-7}$ (proof of Lemma~\ref{lem:oscG}). For fixed $u$, the map
$(x,z_1,z_2)\mapsto G(x,u,z_1,z_2)$ is Lipschitz with constant $\le4(T+1)^2$ except where $u'=z_i'$ (with $u<z_i$) or where the case $u\ge z_i$ changes at
$z_i=0$; the latter does not happen on a clean cell, and the set of $u$ for which $u'-z_i'$ changes sign when the parameters move by $\eta$ has
Lebesgue measure $\le4\eta(T+1)^2$, while the density of $\nu_{x_k,p^C_k}$ is at most $2$ (Corollary~\ref{cor:forget}). This gives the bound on clean cells; on
other cells the same argument applies on each side of the (at most two) crossings of $E_0$, and the values on the two sides differ by at most $10$. The probability bound is
the first part of the proof of \eqref{eq:PB}.
\end{proof}

\begin{lemma}\label{lem:tail}
$\E|\Lambda_{\rm tail}|=o(L)$.
\end{lemma}

\begin{proof}
$|X_k-\E[X_k\mid\cH_J]|\le\osc_{\cC_J(\alpha)}X_k$; apply Lemma~\ref{lem:oscG} with $i=J$ and sum over $k\le M$:
$\E|\Lambda_{\rm tail}|\le M(6e^4L^{-8}+L^{-5})+90\sum_k\Prob(\chi_k=1,\cB_k(\cC_J(\alpha)))=O(L\cdot sL^{-2.7})$.
\end{proof}

\begin{lemma}\label{lem:mart}
$\E[\Lambda_{\rm mart}^2]\le C\eps L^2+o(L^2)$.
\end{lemma}

\begin{proof}
The summands $\ind\{i<J\}D_i$ are $\cH_{i+1}$-measurable with $\E[\ind\{i<J\}D_i\mid\cH_i]=0$, hence orthogonal, and
$\E[\Lambda_{\rm mart}^2]=\sum_i\E[\ind\{i<J\}D_i^2]$. We bound $\|\ind\{i<J\}\Delta_{k,i}\|_2$, recalling that $\Delta_{k,i}=0$ where $\chi_k=0$.
(a) If $i\ge k+3$, $|\Delta_{k,i}|\le\osc_{\cC_i(\alpha)}X_k$, and by Lemma~\ref{lem:oscG},
$\|\ind\{i<J\}\Delta_{k,i}\|_2\le6e^42^{-\lfloor(i-k-2)/2\rfloor}+L^{-5}+90\big(C2^{-\lfloor(i-k-2)/2\rfloor}+C_Q(s+2)L^{-2.7}\big)^{1/2}$.
(b) If $k-1\le i\le k+2$, $|\Delta_{k,i}|\le2(a_{k+1}+5)\ind\{a_{k+1}\le T\}$, whose $L^2$ norm is at most $C\sqrt T$ because $\Prob(a_{k+1}=b)\le2/b^2$.
(c) If $k-s\le i\le k-2$, the cells $C=\cC_i(\alpha)$ and $C'=\cC_{i+1}(\alpha)$ are $k$-regular with $E_C,E_{C'}\le2L^{-2.7}$. Lemma~\ref{lem:Gbarapprox}
applied to both (common reference point $\alpha$), \eqref{eq:Gbarp}, Corollary~\ref{cor:forget} and $(a_{k+1}+5)\alpha_k\le6$ give
\[
|\Delta_{k,i}|\le12\cdot2^{-\lfloor(k-i-2)/2\rfloor}+C(T+5)(s+2)L^{-2.7}+40\cdot\ind\{\cC_{k-s}(\alpha)\text{ is not $k$-clean}\}.
\]
Summing, $\|\ind\{i<J\}D_i\|_2\le C\sqrt T+C+o(1)$, and there are at most $M\le4L$ indices $i$.
\end{proof}

\subsection{The residual}
For $\chi_k=1$, Lemma~\ref{lem:Gbarapprox} with $C=\cC_{k-s}(\alpha)$ gives
\[
\E[X_k\mid\cH_{k-s}]=(-1)^k\chi_k\bar G(\alpha_k,p^C_k,z_{1,k},z_{2,k})+e_k,\qquad \textstyle\sum_k\E|e_k|=o(L).
\] Let $p^\natural_k$ be the solution of \eqref{eq:rec}
along the orbit of $\alpha$ ($x_l=\alpha_l$) started at level $k-s-1$ from $p^\natural_{k-s-1}=\alpha_{k-s-1}$. By Lemma~\ref{lem:contr}(iv),
$|p^C_k-p^\natural_k|\le\alpha_k2^{-\lfloor(s-1)/2\rfloor}$, so by \eqref{eq:Gbarp} replacing $p^C_k$ by $p^\natural_k$ changes each term by at most
$12\cdot2^{-\lfloor(s-1)/2\rfloor}$, in total $O(L2^{-s/2})=o(L)$.

Next, if $p\le2x$ then $\|\nu_{x,p}-\nu_{x,x}\|_{TV}\le x$, so with $\bar\tau(z)=\int_0^1\tau(u,z)\,du=\frac{z(1-z)}2$ and $a=\lfloor1/x\rfloor$,
\begin{equation}\label{eq:Gbarmain}
\big|\bar G(x,p,z_1,z_2)-a\big(\bar\tau(z_2)-\bar\tau(z_1)\big)\big|\le\tfrac12+\tfrac14+4\le5 .
\end{equation}
Moreover $\bar\tau$ is $\frac12$-Lipschitz on the circle and $\|z_{i,k}-u_{i,k}\|\le|v_{i,k}|\alpha_k<4\alpha_k$ (Lemma~\ref{lem:orbit}), so
$a_{k+1}|\bar\tau(z_{i,k})-\bar\tau(u_{i,k})|\le2$. Since $p^\natural_k\le2\alpha_k$ (Lemma~\ref{lem:contr}(ii)), we may write
\[
\ind\{a_{k+1}\le T\}\bar G(\alpha_k,p^\natural_k,z_{1,k},z_{2,k})=f_k+h_k,\qquad f_k=\ind\{a_{k+1}\le T\}\,a_{k+1}\,g(w_k),\quad |h_k|\le9,
\]
with $g(w)=\bar\tau(\fr{P_2w_2/Q})-\bar\tau(\fr{P_1w_2/Q})$. As in Lemma~\ref{lem:R}, on our event
\begin{gather*}
\Big|\Lambda_{\rm res}-\Sigma_1-\Sigma_2\Big|\le\sum_{k\in W'}\ind\{a_{k+1}\le T\}(a_{k+1}+9)+\sum_k|e_k|+o(L),\\
\Sigma_1=\sum_{s<k\le K_0}(-1)^kf_k,\qquad\Sigma_2=\sum_{s<k\le K_0}(-1)^kh_k ,
\end{gather*}
and the window term has expectation $O((A\sqrt L+s)\log L)=o(L)$.

\begin{lemma}\label{lem:digitsums}
$\E[\Sigma_1^2]\le C_Q(\eps L^2+L\log^2L)$, and $\Sigma_2=o_\Prob(L)$; more precisely $\Sigma_2=\Sigma_2'+\Sigma_2''$ with $\E[\Sigma_2'^2]\le C_QL\log L$ and
$|\Sigma_2''|\le L^{-1}$ on $G_N$.
\end{lemma}

\begin{proof}
\emph{$\Sigma_1$.} By (CF6), $\E[f_k\mid\cF_k]=m_T(\omega_k)g(w_k)$ with $m_T(\omega)=\sum_{b\le T}b\,p_\omega(b)\le2(1+\log T)$ and
$|m_T'(\omega)|\le3(1+\log T)$. The martingale differences $f_k-\E[f_k\mid\cF_k]$ are orthogonal, and $\sum_k\E f_k^2\le K_0\cdot2T\le C\eps L^2$. The function
$(\omega,w)\mapsto m_T(\omega)g(w)/(2+2\log T)$ is a $2$-regular functional of the chain at time $k$ ($w_k$ is a function of $\zeta_k$). By
Lemma~\ref{lem:mixing}(b) its means converge exponentially fast, so $|\sum_k(-1)^k\E[m_T(\omega_k)g(w_k)]|\le C_Q\log T$, and by Lemma~\ref{lem:mixing}(a)
its covariances at lag $n$ are $O_Q(\vartheta_Q'^{\,n}\log^2T)$. Hence $\E[(\sum_k(-1)^km_T(\omega_k)g(w_k))^2]\le C_QL\log^2L$.

\emph{$\Sigma_2$.} The variable $h_k$ is determined by $\zeta_{k-s-1}$ and the pieces $\fp_{k-s-1},\fp_{k-s},\dots$: the pieces determine the digits and the
$\alpha_l$, $p^\natural_k$ is computed from $\alpha_{k-s-1},\dots,\alpha_k$ and the digits, and $z_{i,k}=\fr{u_{i,k}+v_{i,k}\alpha_k}$ with $\zeta_k$ obtained from
$\zeta_{k-s-1}$ by \eqref{eq:zetarec}. So $h_k=\phi(\zeta_{k-s-1},\fp_{k-s-1},\fp_{k-s},\dots)$ for a Borel function $\phi$ (depending on $N$ only through
$s,T$) with $|\phi|\le9$, independent of $\omega$. By Lemma~\ref{lem:mixing}(b) the alternating sum of the means $\E h_k$ is $O_Q(1)$. Let $r=\lceil20\log L\rceil$ and
let $h^{[r]}_k$ be obtained by replacing every $\alpha_l$ ($k-s-1\le l\le k+1$) by $\alpha_l^{[r+k+1-l]}$ in the formula for $h_k$ and then clipping the
result to $[-9,9]$ (the substituted values need not form a Gauss orbit, so the bound $|h_k|\le9$ is not automatic); this is a function of $\zeta_{k-s-1}$ and
$\fp_{k-s-1},\dots,\fp_{k+r+1}$, bounded by $9$. Inside fixed pieces, $\phi$ is Lipschitz in these $\alpha_l$ with constant at most $L^{10}$ on $G_N$: the maps
$\lambda(\cdot),\mu(\cdot)$ of \eqref{eq:lambdamu} have derivative $O(x^{-2})=O(U^2)$, $|\partial_p\bar G|\le2(T+5)$, the density of $\nu_{x,p}$ depends on
$x$ with derivative $O(T\,U)$, the integrals of the $\tilde\kappa$-terms over $u$ are Lipschitz in $(x,z_i)$ with constant $O((T+1)^2)$ (the jumps of
$\tilde\kappa$ in $u$ only move), and $z_{i,l}$ is affine in $\alpha_l$ inside a piece. Since clipping is $1$-Lipschitz and $|h_k|\le9$, $|h_k-h_k^{[r]}|\le L^{10}F_{r+1}^{-2}\le L^{-5}$ on $G_N$. Put
$\Sigma_2'=\sum_k(-1)^kh^{[r]}_k$ and $\Sigma_2''=\Sigma_2-\Sigma_2'$. The alternating sum of the means of $h^{[r]}_k$ is $O_Q(1)$ (Lemma~\ref{lem:mixing}(b)), and for
$k-j>s+r+3$, Lemma~\ref{lem:mixing}(a) gives $|\mathrm{Cov}(h^{[r]}_j,h^{[r]}_k)|\le C_Q\vartheta_Q'^{\,k-j-s-r-3}$. Hence
$\E[\Sigma_2'^2]\le C_QL(s+r)$.
\end{proof}

\begin{proof}[Proof of Proposition~\ref{prop:light}]
Combine $\Lambda=\Lambda_{\rm res}+\Lambda_{\rm mart}+\Lambda_{\rm tail}$ with Lemmas~\ref{lem:good}, \ref{lem:bulkend}, \ref{lem:sigmaZ}, \ref{lem:tail}, \ref{lem:mart},
\ref{lem:digitsums} and Chebyshev's inequality; $\eta$ is arbitrary.
\end{proof}
\section{The large partial quotients}\label{sec:heavy}

For $c\in\Z/Q\Z$ let $m(t,c)=\tau(t,\fr{P_2c/Q})-\tau(t,\fr{P_1c/Q})$ (with the convention $\tau(t,\fr{P_ic/Q})=0$ for a zero endpoint), so that
$|m|\le\frac14$, and let
\[
\varphi_\eps(\theta)=\exp\Big(\frac{c_0}{\log2}\,\E\int_\eps^\infty\!\!\int_0^1\big(\cos(\theta x\,m(t,c_Q))-1\big)\,dt\,\frac{dx}{x^2}\Big),
\]
the characteristic function of a symmetric compound Poisson law $\Pi_\eps$.

\begin{proposition}\label{prop:heavy}
$H/L$ converges in distribution to $\Pi_\eps$ as $N\to\infty$.
\end{proposition}

\subsection*{Reductions}
Put $m_k(t)=m(t,(w_k)_2)$; thus $m_k(t)=\tau(t,u_{2,k})-\tau(t,u_{1,k})$, and $m_k$ is $\cF_k$-measurable. By \eqref{eq:LMZK} and
$\|z_{i,k}-u_{i,k}\|<4\alpha_k$,
\[
|L_k-a_{k+1}m_k(y_k)|\le\alpha_{k+1}\cdot\tfrac14+a_{k+1}\sum_i\|z_{i,k}-u_{i,k}\|+|Z_k|+|K_k|\le14 .
\]
Since $\E\#\{k\le M:a_{k+1}>T\}\le2M/T\le7/\eps$, we have $|H-\sum_k(-1)^k\chi'_ka_{k+1}m_k(y_k)|\le14\#\{k\le M:a_{k+1}>T\}$. Fix $\eta>0$, $A=A(\eta)$ as in
Lemma~\ref{lem:bulkend}, put $K_*=K_0-\lceil A\sqrt L\rceil$ and
\[
\xi_k=\ind\{T<a_{k+1}\le U\}\,\ind\{s<k\le K_*\},\qquad H_2=\sum_k(-1)^k\xi_ka_{k+1}m_k(y_k).
\]
On $G_N\cap\{|K''-K_0|\le A\sqrt L\}$ we have $\chi'_k=\ind\{a_{k+1}>T\}\ind\{s<k\le K''+s\}$, so $\chi'_k\ne\xi_k$ forces a partial quotient
$a_{k+1}>T$ with $|k-K_0|\le A\sqrt L+s+1$, an event of probability $O(\sqrt L/T)$. Hence, off an event of probability $\eta+o(1)$,
$|H-H_2|\le14\#\{k\le M:a_{k+1}>T\}$, whose expectation is $O(1/\eps)$. Consequently
\begin{equation}\label{eq:Hred}
\limsup_{N\to\infty}\big|\E e^{i\theta H/L}-\varphi_\eps(\theta)\big|\le2\eta+\limsup_{N\to\infty}\big|\E e^{i\theta H_2/L}-\varphi_\eps(\theta)\big|,
\end{equation}
and it suffices to show that the last $\limsup$ is at most $C(\theta)\sqrt\eta/\eps$; since the left side of \eqref{eq:Hred} does not depend on
$\eta$, letting $\eta\to0$ then proves Proposition~\ref{prop:heavy}. Put
\[
\Theta_k=e^{i\theta(-1)^ka_{k+1}m_k(y_k)/L}-1,\qquad \hat\Theta_k=\int_0^1\big(e^{i\theta(-1)^ka_{k+1}m_k(t)/L}-1\big)\,dt,
\]
$A_k=1+\xi_k\Theta_k$ and $B_k=1+\xi_k\hat\Theta_k$; note $|A_k|,|B_k|\le1$, $e^{i\theta H_2/L}=\prod_kA_k$, and $B_k$ is $\cF_{k+1}$-measurable.

\begin{lemma}[Poisson limit for the digits and residues]\label{lem:poisson}
$\E\prod_kB_k\to\varphi_\eps(\theta)$.
\end{lemma}

\begin{proof}
Let $\Gamma_k=\E[B_k\mid\cF_k]$. By (CF6), $|\Gamma_k-1|\le2\Prob(a_{k+1}>T\mid\cF_k)\le4/T$, so $|\Gamma_k|\ge\frac12$ and
$\prod_k|\Gamma_k|^{-1}\le e^{8M/T}\le e^{24/\eps}$. Since $\E[B_k/\Gamma_k\mid\cF_k]=1$, the product $\prod_{k\le n}B_k/\Gamma_k$ has expectation $1$ for every $n$, so
\[
\Big|\E\prod_kB_k-\varphi_\eps(\theta)\Big|=\Big|\E\Big[\prod_k\frac{B_k}{\Gamma_k}\Big(\prod_k\Gamma_k-\varphi_\eps(\theta)\Big)\Big]\Big|\le e^{24/\eps}\,\E\Big|\prod_k\Gamma_k-\varphi_\eps(\theta)\Big| ,
\]
and it suffices to show $\prod_k\Gamma_k\to\varphi_\eps(\theta)$ in probability (all quantities being bounded). As $\sum_k|\Gamma_k-1|^2\le16M/T^2\to0$, it suffices
to show $\sum_k(\Gamma_k-1)\to\log\varphi_\eps(\theta)$ in probability. By (CF6),
\[
\Gamma_k-1=\xi^{(0)}_k\sum_{T<b\le U}p_{\omega_k}(b)\int_0^1\big(e^{i\theta(-1)^kbm_k(t)/L}-1\big)dt,\qquad \xi^{(0)}_k=\ind\{s<k\le K_*\},
\]
and replacing $p_{\omega_k}(b)$ by $(1+\omega_k)b^{-2}$ costs $\sum_{b>T}20b^{-3}\le20T^{-2}$ per $k$. A Riemann sum with $x=b/L$ (the integrand has
derivative $O((|\theta|+\eps^{-1})x^{-2})$ on $[\eps,\infty)$) and the bound $\int_{U/L}^\infty2x^{-2}dx=2L/U$ give
\[
\sum_k(\Gamma_k-1)=\frac1L\sum_{s<k\le K_*}(1+\omega_k)\Big(R(c_k)+i(-1)^kI(c_k)\Big)+o(1),
\]
where $c_k=(w_k)_2$ and, for $c\in\Z/Q\Z$,
\[
R(c)=\int_\eps^\infty\!\!\int_0^1\big(\cos(\theta xm(t,c))-1\big)dt\,\frac{dx}{x^2},\qquad I(c)=\int_\eps^\infty\!\!\int_0^1\sin(\theta xm(t,c))\,dt\,\frac{dx}{x^2};
\]
note $|R|,|I|\le2/\eps$.
The functionals $(1+\omega_k)R(c_k)\eps/4$ and $(1+\omega_k)I(c_k)\eps/4$ of the chain at time $k$ are $1$-regular. By Lemma~\ref{lem:statlaw},
$\E[(1+\omega_k)R(c_k)]=\frac1{\log2}\E R(c_Q)+O(\vartheta_Q'^{\,k}/\eps)$, because $\int(1+t)\gamma(dt)=1/\log2$; and by Lemma~\ref{lem:mixing}(a) the covariances
of these functionals at lag $n$ are $O(\vartheta_Q'^{\,n}/\eps^2)$. Hence, in $L^2$,
\[
\frac1L\sum_{s<k\le K_*}(1+\omega_k)R(c_k)\to\frac{c_0}{\log2}\E R(c_Q),\qquad \frac1L\sum_{s<k\le K_*}(-1)^k(1+\omega_k)I(c_k)\to0,
\]
using $(K_*-s)/L\to c_0$. Since $\log\varphi_\eps(\theta)=\frac{c_0}{\log2}\E R(c_Q)$, the lemma follows.
\end{proof}

It remains to show $\E\prod_kA_k-\E\prod_kB_k\to0$. By telescoping,
\begin{equation}\label{eq:telescope}
\Big|\E\prod_kA_k-\E\prod_kB_k\Big|\le\sum_{k'}\Big|\E\Big[\prod_{k<k'}A_k\cdot\xi_{k'}(\Theta_{k'}-\hat\Theta_{k'})\cdot\prod_{k>k'}B_k\Big]\Big|.
\end{equation}
The difficulty is that $y_{k'}$ is \emph{not} independent of the later digits that enter $\prod_{k>k'}B_k$: the cells reveal digits
far ahead. We separate the later indices into those which are already determined by the cell $\cC_{k'-s}(\alpha)$ and those which are still
``mixing'' given a cylinder on which $y_{k'}$ is essentially constant.

\subsection*{Near, far and transition indices}
Fix $k'$ with $\xi_{k'}=1$ and put $X=Nq_{k'}$. Let
\[
t_n=\max\{t\ge k':\ q_{t+2}^2\le L^{-6}X\},\qquad t_f=\min\{t:\ q_t^2\ge L^4X\},
\]
\[
r_h=\lceil\log^2L\rceil,\qquad w'=\lceil10\log_2L\rceil+6+r_h,
\]
and call $k>k'$ \emph{near} if $k\le t_n+1$, \emph{transition} if $t_n+1<k<t_n+w'$, and \emph{far} if $k\ge t_n+w'$. Here
$\{t_n=t\}\in\sigma(a_1,\dots,a_{t+3})$ and $t_f$ is a stopping time of the digit filtration. Since $q_{t+2}\ge2q_t$, the definitions give
\begin{equation}\label{eq:zone}
t_n+3\le t_f\le t_n+w'-r_h .
\end{equation}
Write $\prod_{k>k'}B_k=F_{\rm near}F_{\rm tr}F_{\rm far}$ accordingly, and $\prod_{k<k'}A_k=P'P_{\rm np}$ where $P_{\rm np}$ collects the
indices $k'-2s-2\le k<k'$.

\begin{lemma}[Exceptional events]\label{lem:exc}
Let $\mathcal E_{k'}$ be the union of the events: (i) $G_N^c$, or $|K''-K_0|>A\sqrt L$, or $g_{k'}=0$; (ii) $\xi_k=1$ for some $k\ne k'$ with
$|k-k'|\le2s+2$; (iii) the set defining $t_n$ is empty, or $\xi_k=1$ for some transition index $k$; (iv) $\cC_{k'-s}(\alpha)\not\subset D_{t_n+3}(\alpha)$,
the cylinder of depth $t_n+3$ containing $\alpha$; (v) $D_{t_f}(\alpha)\not\subset\cC_{k'+1}(\alpha)$. Then
$\sum_{k'}\E[\xi_{k'}\ind_{\mathcal E_{k'}}]\le C\sqrt\eta/\eps+o(1)$.
\end{lemma}

\begin{proof}
This involves the digits and the cells only, and the proof of \cite[Lemma~8.3]{E1002} applies verbatim; we recall it.
(i) On $G_N\cap\{|K''-K_0|\le A\sqrt L\}$ every $k'\le K_*$ has $g_{k'}=1$; and $\E[\#\{k:\xi_k=1\}^2]\le C/\eps^2$ by (CF2), so the contribution is
at most $(C/\eps)(\eta+o(1))^{1/2}$.
(ii) $\sum_{k'}\Prob(\xi_{k'}=\xi_k=1,\ 0<|k'-k|\le2s+2)\le CMs/T^2\to0$ by (CF2).
(iii) Assume (i) fails. Then $k'\le K_*\le K''$, so $N_{k'}\ge Z$ and $q_{k'}<N/N_{k'}\le N/Z$ (Lemma~\ref{lem:sizes}). As
$q_{k'+2}\le(a_{k'+1}+1)(a_{k'+2}+1)q_{k'}$ and $a_{k'+1}\le U$, emptiness forces $a_{k'+2}>L^{0.9}/2$, which has probability $O(L^{-0.9})$ given $\xi_{k'}=1$. Transition indices
$k\ge t_n+3$ have $a_{k+1}$ beyond the digits determining $\{t_n=t\}\cap\{\xi_{k'}=1\}$, so by
(CF2) each is heavy with conditional probability $\le C/T$, and there are at most $w'$ of them. For $k=t_n+2$, the event
$\{t_n=t,\ \xi_{t+2}=1\}$ (which entails $T<a_{t+3}\le U$) forces $q_{t+2}^2\in(L^{-6}X(U+1)^{-2},L^{-6}X]$, which happens for at most $4\log_2(U+1)+2$ values of $t$,
each time with conditional probability $\le2/T$ given $\sigma(a_1,\dots,a_{t+2})$. Altogether this gives $C(w'+\log L)/T$ times $\Prob(\xi_{k'}=1)\le2/T$, and the
sum over $k'\le M$ is $O(L\log^2L/T^2)\to0$.
(iv) Assume (i)--(iii) fail. Since $q_{t_n+2}^2\le L^{-6}Nq_{k'}<L^{-16}N^2$, we have $t_n+2<m\le M$ (with $m$ as in Lemma~\ref{lem:sizes}),
so $a_{t_n+3}\le U$ on $G_N$. Further $|\cC_{k'-s}|\le4/(Nq_{k'-s-1})$ (proof of Lemma~\ref{lem:good}),
$\lambda(D_{t_n+3})\ge(2(U+1)^2q_{t_n+2}^2)^{-1}\ge L^6/(2(U+1)^2X)$ and $q_{k'}\le q_{k'-s-1}\Pi_{k'}$, so $|\cC_{k'-s}|\le L^{-2.5}\lambda(D_{t_n+3})$.
Hence (iv) implies $\mathrm{dist}(\alpha,\partial D_{t+3}(\alpha))\le L^{-2.5}\lambda(D_{t+3}(\alpha))$ for $t=t_n$. For fixed $t\ge k'$ the event
$\{a_{k'+1}>T\}$ is a union of cylinders of depth $t+3$, so
$\Prob\big(a_{k'+1}>T,\ \mathrm{dist}(\alpha,\partial D_{t+3}(\alpha))\le\epsilon\lambda(D_{t+3}(\alpha))\big)=2\epsilon\Prob(a_{k'+1}>T)$; summing over
$t\le M$ and $k'\le M$ gives $O(M^2L^{-2.5}/T)\to0$.
(v) Assume (i)--(iv) fail. The range of the level-$(k'+1)$ cell is of type \textbf F, \textbf P or \textbf Q (Proposition~\ref{prop:types}) with
parameter within $\osc_{\cC_{k'-s}}\alpha_{k'}\le4\Pi_{k'}^2/N_{k'-s-1}\le4L^{-7.6}$ of $\alpha_{k'}$ (Lemma~\ref{lem:osc}(b)), and
$\alpha_{k'},1-\alpha_{k'}\ge(2(U+1))^{-1}$, so its length is $\ge(4(U+1))^{-1}$; by the formula $|d\alpha/dy_{k'}|\ge(4N_{k'}q_{k'}^2)^{-1}$ from the
proof of Lemma~\ref{lem:distortion}, $|\cC_{k'+1}|\ge(16(U+1)N_{k'}q_{k'}^2)^{-1}\ge(16(U+1)X)^{-1}$, whereas $\lambda(D_{t_f})\le q_{t_f}^{-2}\le L^{-4}X^{-1}$. Thus (v) implies
$\mathrm{dist}(\alpha,\partial\cC_{k'+1}(\alpha))\le32L^{-2.9}|\cC_{k'+1}(\alpha)|$. The level-$(k'+1)$ cells partition $(0,1)$ and $\{a_{k'+1}>T\}$ is a
union of them, so this has probability $\le64L^{-2.9}\Prob(a_{k'+1}>T)$; sum over $k'$.
\end{proof}

\begin{lemma}\label{lem:term}
For each $k'$, the $k'$-th term of \eqref{eq:telescope} is at most
$C\E[\xi_{k'}\ind_{\mathcal E_{k'}}]+C\,\Prob(\xi_{k'}=1)\big(T^{-1}+sL^{-2.7}+\vartheta_Q'^{\,r_h}+|\theta|UL^{-5}\big)+C|\theta|2^{-s/2}L^{-1}\E[\xi_{k'}\#\{k:\xi_k=1\}]$.
\end{lemma}

\begin{proof}
Off $\mathcal E_{k'}$ we have $g_{k'}=1$, $P_{\rm np}=1$, $F_{\rm tr}=1$ and the set defining $t_n$ is non-empty; inserting the factor
$g_{k'}\ind\{\text{the set defining }t_n\text{ is non-empty}\}$ and replacing $P_{\rm np}$ and $F_{\rm tr}$ by $1$ costs at most
$C\E[\xi_{k'}\ind_{\mathcal E_{k'}}]$. The non-emptiness indicator depends only on $q_{k'},q_{k'+2}$, hence on $a_1,\dots,a_{k'+2}$, and is
therefore constant on $k'$-regular level-$(k'-s)$ cells. Put $\Delta=\Theta_{k'}-\hat\Theta_{k'}$.

\emph{The past.} On $\{g_{k'}=1\}$ we have $J>k'+1$, so Lemma~\ref{lem:osc}(a) applies on $\cC_{k'-s}(\alpha)$. Each factor of $P'$ with $\xi_k=1$ involves
$y_k$ with $k\le k'-2s-3$ and the function $m_k$, which is constant on $\cC_{k'-s}(\alpha)$; as $|m_k(t)-m_k(t')|\le2|t-t'|$,
$\osc_{\cC_{k'-s}}\Theta_k\le2|\theta|a_{k+1}\osc y_k/L\le C|\theta|2^{-s/2}/L$. Let $P''$ be the function
obtained from $P'$ by freezing every $y_k$ ($k<k'-2s-2$) at its value at a fixed point of the level-$(k'-s)$ cell of $\alpha$; then
$|P''|\le1$, $P''$ is constant on level-$(k'-s)$ cells, and $|P'-P''|\le C|\theta|2^{-s/2}L^{-1}\#\{k:\xi_k=1\}$.

\emph{The near future.} If (iv) fails, $\cC_{k'-s}(\alpha)$ lies in a single cylinder of depth $t_n+3$; then $t_n$ and $F_{\rm near}$ are constant
on the cell (the factors $B_k$, $k\le t_n+1$, depend on $a_{k+1}$ and $w_k$, i.e.\ on $a_1,\dots,a_{t_n+2}$), and so is the indicator of (iv)$^c$. Thus
$Y:=P''\xi_{k'}g_{k'}\ind\{\text{non-empty}\}\ind_{\rm (iv)^c}F_{\rm near}$ is $\cH_{k'-s}$-measurable and $|Y|\le1$.

\emph{The far future.} On $\{t_n=t\}$, $F_{\rm far}=\Phi_t$, where $\Phi_t=\prod_{k\ge t+w'}B_k$ is a functional of the extended chain from time $t+w'$ (it
depends on $\zeta_{t+w'}$, through $w_k$, and on the subsequent pieces), bounded by $1$ and independent of $\omega$. If (v) fails,
$D:=D_{t_f}(\alpha)\subset\cC_{k'+1}(\alpha)$; there $|dy_{k'}/d\alpha|\le4N_{k'}q_{k'}^2\le4X$, so $\osc_Dy_{k'}\le4L^{-4}$ and
$\osc_D\Delta\le8|\theta|UL^{-5}$; moreover $Y$, $t_n$ and $\ind_{(v)^c}$ are constant on $D$ (note $t_n+3\le t_f$ and $D\subset\cC_{k'-s}$). Since $D$ is a cylinder of
depth $t_f\le t_n+w'-r_h$ \eqref{eq:zone}, Lemma~\ref{lem:mixing}(a) gives $\E[F_{\rm far}\mid D]=\phi_{t_n}+O(\vartheta_Q'^{\,r_h})$ with $\phi_t=\E\Phi_t$. Hence
\[
\E\big[Y\Delta\ind_{(v)^c}F_{\rm far}\big]=\E\big[Y\phi_{t_n}\Delta\ind_{(v)^c}\big]+O\big(\Prob(\xi_{k'}=1)(\vartheta_Q'^{\,r_h}+|\theta|UL^{-5})\big).
\]
Removing $\ind_{(v)^c}$ costs $C\E[\xi_{k'}\ind_{\mathcal E_{k'}}]$, and $Y\phi_{t_n}$ is $\cH_{k'-s}$-measurable, so we are left with
$\E\big[Y\phi_{t_n}\E[\Delta\mid\cH_{k'-s}]\big]$. On $\{\xi_{k'}g_{k'}=1\}$, $\Theta_{k'}$ is a function of $y_{k'}$ bounded by $2$ ($m_{k'}$ being constant on the cell), so
Theorem~\ref{thm:condlaw} and Lemma~\ref{lem:contr}(ii) give, with $p=p^{\cC_{k'-s}}_{k'}\le2\alpha_{k'}$,
\[
|\E[\Delta\mid\cH_{k'-s}]|\le4\big\|\mathrm{Law}(y_{k'}\mid\cC_{k'-s})-\nu_{\alpha_{k'},p}\big\|_{TV}+4\big\|\nu_{\alpha_{k'},p}-\nu_{\alpha_{k'},\alpha_{k'}}\big\|_{TV}
\le C(s+2)L^{-2.7}+8\alpha_{k'},
\]
with $\alpha_{k'}<1/T$.
\end{proof}

\begin{proof}[Proof of Proposition~\ref{prop:heavy}]
Sum Lemma~\ref{lem:term} over $k'$. Using $\sum_{k'}\Prob(\xi_{k'}=1)\le2M/T\le7/\eps$, $\E[\#\{k:\xi_k=1\}^2]\le C/\eps^2$ and Lemma~\ref{lem:exc}, the
right-hand side of \eqref{eq:telescope} is at most $C(\theta)\sqrt\eta/\eps+o(1)$. With Lemma~\ref{lem:poisson} this gives
$\limsup_N|\E e^{i\theta H_2/L}-\varphi_\eps(\theta)|\le C(\theta)\sqrt\eta/\eps$, and \eqref{eq:Hred} with $\eta\to0$ gives
$\E e^{i\theta H/L}\to\varphi_\eps(\theta)$ for every $\theta$.
\end{proof}
\section{Proof of Theorem~\ref{thm:main}}\label{sec:end}

By \eqref{eq:decomp}, $D_N(\alpha,\beta;r)/L=(\Lambda+\Lambda_Z)/L+H/L+R/L$. By Lemma~\ref{lem:R}, $R/L\to0$ in probability. For fixed $\eps$,
$H/L\Rightarrow\Pi_\eps$ (Proposition~\ref{prop:heavy}) and $\limsup_N\Prob(|\Lambda+\Lambda_Z|/L>\delta)\le C_Q\eps/\delta^2$ (Proposition~\ref{prop:light}).
Since $\int_0^\infty(1-\cos cx)x^{-2}dx=\frac\pi2|c|$ and $\int_0^\eps(1-\cos cx)x^{-2}dx\le c^2\eps/2$,
\[
\lim_{\eps\to0}\log\varphi_\eps(\theta)=-\frac{c_0}{\log2}\cdot\frac{\pi|\theta|}{2}\cdot\E\int_0^1|m(t,c_Q)|\,dt=-\frac6\pi\,\E F\big(\fr{z_1c_Q},\fr{z_2c_Q}\big)\,|\theta|=-\sigma(\beta,r)|\theta| ,
\]
because $\frac{c_0}{\log2}\cdot\frac\pi2=\frac6\pi$ and $\int_0^1|m(t,c)|dt=F(\fr{P_1c/Q},\fr{P_2c/Q})$ with $P_i/Q=z_i$. So $\Pi_\eps$ converges to the Cauchy law with
characteristic function $e^{-\sigma(\beta,r)|\theta|}$. By the standard approximation theorem \cite[Theorem~3.2]{Bi99}, $D_N(\alpha,\beta;r)/\log N$ converges in
distribution to this law when $\alpha\sim\lambda$. Its distribution function $\frac12+\frac1\pi\arctan(c/\sigma(\beta,r))$ is continuous, which proves
Theorem~\ref{thm:main}. \qed

\begin{remark}[Where rationality is used]\label{rem:where}
Rationality of $\beta$ and $r$ enters only through Section~\ref{sec:orbits}: Lemma~\ref{lem:orbit} (the endpoint orbits live on finitely many lines
$z=\fr{u+v\alpha_k}$), the Markov refinement at the finite set $E_0$, and the coupling Lemma~\ref{lem:coupling}. For irrational endpoints the same
decomposition holds, and the proof would go through verbatim given a law of large numbers, with a rate $o(1/\log L)$ for the alternating version, for
functionals of $(\omega_k,\fr{z_iq_k})$; this is the equidistribution problem for the orbit $(\fr{z_iq_{k-1}},\fr{z_iq_k})$ of the torus under the random
products of $\left(\begin{smallmatrix}0&1\\1&a\end{smallmatrix}\right)$ driven by the continued fraction digits, which we do not address. Borda raised the
analogous question for the maximal sums \cite{Bo25}.
\end{remark}

\begin{remark}[Symmetries]
For irrational $\alpha$ one checks $D_N(1-\alpha,\beta;r)=D_N(\alpha,r-\beta;r)$ and $D_N(\alpha,r;r)=-D_N(\alpha,0;1-r)$. For $\beta=0$, $r=\frac12$ this gives
$D_N(1-\alpha;\frac12)=-D_N(\alpha;\frac12)$, so in that case the symmetry of the limit law is obvious. In general these identities relate different
starting points, and the centring of the limit law comes from the parity-free dynamics of Section~\ref{sec:orbits}; consistently with them,
$\sigma(\beta,r)=\sigma(r-\beta,r)$ and $\sigma(r,r)=\sigma(0,1-r)$, since $\fr{-\beta}$ and $\fr{r-\beta}$ are then exchanged or reflected.
\end{remark}
\section{Numerical checks}\label{sec:numerics}

All computations use exact integer arithmetic. We take $\alpha=A/2^B$ with $A$ uniform and $2^B\ge N^{4}$, evaluate $D_N(\alpha;z)$ by a
Euclid-type floor-sum algorithm, and use $D_N(\alpha,\beta;r)=D_N(\alpha;z_2)-D_N(\alpha;z_1)$ (Lemma~\ref{lem:reduction}). Everything below is
reproduced by the script \texttt{verify\_kesten.py} accompanying this paper.

\smallskip\noindent\textbf{Exact identities.} (a) The expansion of Proposition~\ref{prop:expansion} was compared with the direct count for $600$ random
triples $(\alpha,z,N)$ ($\alpha$ with denominators $2^{96}$--$2^{256}$, $z=P/Q$ with $Q\le1000$, $N$ up to $10^{14}$): no mismatch.
(b) The structure of Lemma~\ref{lem:orbit} ($z_k=\fr{u_k+v_k\alpha_k}$, $v_k\equiv P(w_k)_1/Q$, $-2-2\alpha_{k-1}<v_k\le0$) was checked at $28000$ levels
($Q\in\{2,3,5,7,12\}$): no violation; the largest $|v_k|$ observed was $2.33$, and the reset $v_{k+1}=-u_k$ after a partial quotient $\ge4Q+1$ with $u_k\ne0$
(Step~2 of Lemma~\ref{lem:coupling}) held in all $1786$ instances. (c) For $y$ uniform, $\int_0^1\ell(x,y,z)\,dy=\frac{z(1-z)}{2x}+z'+\frac z2-1-\frac{xz'(1-z')}2$
with $z'=\fr{-z/x}$; this was verified exactly for $300$ rational pairs $(x,z)$. It shows that the bounded part of the level function has a non-zero
mean, depending on the endpoint orbit; its alternating sum is $o(L)$ only because even and odd levels have the same statistics (Section~\ref{sec:orbits}).

\smallskip\noindent\textbf{The limit law at rational parameters.} For $N=10^{30}$ and $20000$ independent $\alpha$ per case we list quantiles of
$X=D_N(\alpha,\beta;r)/\log N$ and the scale estimates $\hat\sigma(\theta)=-\theta^{-1}\log|\frac1n\sum\cos(\theta X)|$ (for a Cauchy law with scale $\sigma$,
$\hat\sigma(\theta)\to\sigma$), next to the exact value of $\sigma(\beta,r)$:
\begin{center}\small
\begin{tabular}{lcccccccc}\toprule
$(\beta,r)$ & $\pi\sigma(\beta,r)$ & $\sigma(\beta,r)$ & $q_{.10}$ & $q_{.25}$ & $q_{.75}$ & $q_{.90}$ & $\hat\sigma(1)$ & $\hat\sigma(2)$\\\midrule
$(0,\frac12)$ & $\frac12$ & $.1592$ & $-.478$ & $-.159$ & $.159$ & $.478$ & $.1551$ & $.1573$\\
$(0,\frac13)$ & $\frac12$ & $.1592$ & $-.483$ & $-.150$ & $.154$ & $.473$ & $.1543$ & $.1569$\\
$(0,\frac27)$ & $\frac12$ & $.1592$ & $-.467$ & $-.163$ & $.155$ & $.474$ & $.1524$ & $.1576$\\
$(\frac14,\frac12)$ & $\frac14$ & $.0796$ & $-.246$ & $-.087$ & $.072$ & $.232$ & $.0782$ & $.0777$\\
$(\frac15,\frac25)$ & $\frac14$ & $.0796$ & $-.246$ & $-.087$ & $.087$ & $.232$ & $.0765$ & $.0786$\\
$(\frac23,\frac13)$ & $\frac14$ & $.0796$ & $-.236$ & $-.077$ & $.082$ & $.256$ & $.0744$ & $.0785$\\
$(\frac17,\frac37)$ & $\frac3{10}$ & $.0955$ & $-.296$ & $-.108$ & $.095$ & $.298$ & $.0941$ & $.0965$\\
$(\frac13,\frac12)$ & $\frac13$ & $.1061$ & $-.333$ & $-.116$ & $.116$ & $.333$ & $.1091$ & $.1074$\\
$(\frac17,\frac12)$ & $\frac13$ & $.1061$ & $-.318$ & $-.116$ & $.101$ & $.318$ & $.1004$ & $.1035$\\
$(\frac7{10},\frac3{10})$ & $\frac{499}{1512}$ & $.1051$ & $-.333$ & $-.116$ & $.116$ & $.318$ & $.1030$ & $.1041$\\
$(\frac12,\frac13)$ & $\frac38$ & $.1194$ & $-.352$ & $-.121$ & $.125$ & $.357$ & $.1129$ & $.1163$\\\bottomrule
\end{tabular}
\end{center}
For a Cauchy law with scale $\sigma$ the quartiles are $\pm\sigma$ and the deciles $\pm3.078\,\sigma$. The quantiles are affected by the lattice
structure of $D_N$ (its values are spaced by $1/\log N\approx0.0145$ here), the estimates $\hat\sigma(\theta)$ are not. The estimates $\hat\sigma(2)$
(standard error about $0.0035$) agree with the predicted scales to within $1$--$3\%$, the noisier $\hat\sigma(1)$ (standard error about $0.006$) to
within one or two standard errors, and both separate the different values of $\sigma(\beta,r)$ clearly. A small downward bias of a few per cent
was also observed for the sawtooth sums in \cite{E1002}; it reflects the slow convergence (there are only about $0.84\log N\approx58$
renormalisation steps).

For $N=10^{60}$ ($12000$ samples per case) the estimates are:
\begin{center}\small
\begin{tabular}{lccccccc}\toprule
$(\beta,r)$ & $\sigma(\beta,r)$ & $q_{.10}$ & $q_{.25}$ & $q_{.75}$ & $q_{.90}$ & $\hat\sigma(1)$ & $\hat\sigma(2)$\\\midrule
$(0,\frac13)$ & $.1592$ & $-.458$ & $-.154$ & $.157$ & $.454$ & $.1547$ & $.1528$\\
$(\frac14,\frac12)$ & $.0796$ & $-.246$ & $-.080$ & $.080$ & $.239$ & $.0763$ & $.0782$\\
$(\frac17,\frac37)$ & $.0955$ & $-.293$ & $-.097$ & $.098$ & $.286$ & $.0954$ & $.0966$\\
$(\frac7{10},\frac3{10})$ & $.1051$ & $-.326$ & $-.109$ & $.109$ & $.326$ & $.1074$ & $.1049$\\
$(\frac12,\frac13)$ & $.1194$ & $-.372$ & $-.118$ & $.128$ & $.374$ & $.1182$ & $.1225$\\\bottomrule
\end{tabular}
\end{center}
The standard error of $\hat\sigma(2)$ is about $0.003$ for $20000$ samples and $0.004$ for $12000$ samples, so all entries are compatible with the
predicted scales.

\smallskip\noindent\textbf{Centring.} The most delicate point of the proof is that the drift produced by the bounded terms is $o(\log N)$
(Section~\ref{sec:light}). We estimated the location $\hat c(\theta)=\theta^{-1}\arg\big(\frac1n\sum e^{i\theta X}\big)$ for two asymmetric pairs,
with $40000$ samples per value of $N$:
\begin{center}\small
\begin{tabular}{lcccc}\toprule
$N$ & $10^{12}$ & $10^{20}$ & $10^{45}$ & $10^{60}$\\\midrule
$(\frac17,\frac37)$: $\hat c(2)$ & $-.0094$ & $-.0075$ & $-.0021$ & $-.0019$\\
\phantom{$(\frac17,\frac37)$:} $\hat c(2)\log N$ & $-.26$ & $-.35$ & $-.21$ & $-.27$\\
$(\frac7{10},\frac3{10})$: $\hat c(2)$ & $.0052$ & $.0020$ & $.0028$ & $.0016$\\
\phantom{$(\frac7{10},\frac3{10})$:} $\hat c(2)\log N$ & $.15$ & $.09$ & $.29$ & $.22$\\\bottomrule
\end{tabular}
\end{center}
(standard error of $\hat c(2)$ about $0.0013$). The location of $D_N$ itself thus stays bounded, as predicted by the proof, and $\hat c\to0$.

\smallskip\noindent\textbf{Irrational parameters (Conjecture~\ref{conj:irr}).} With the same $N$ and sample size:
\begin{center}\small
\begin{tabular}{lcccccccc}\toprule
$(\beta,r)$ & conjectured $\sigma$ & $q_{.10}$ & $q_{.25}$ & $q_{.75}$ & $q_{.90}$ & $\hat\sigma(1)$ & $\hat\sigma(2)$\\\midrule
$(0,\sqrt2-1)$ & $\frac1{2\pi}=.1592$ & $-.461$ & $-.157$ & $.162$ & $.466$ & $.1516$ & $.1548$\\
$(0,1/\pi)$ & $\frac1{2\pi}=.1592$ & $-.471$ & $-.152$ & $.152$ & $.470$ & $.1538$ & $.1551$\\
$(0,\frac{\sqrt5-1}2)$ & $\frac1{2\pi}=.1592$ & $-.473$ & $-.155$ & $.149$ & $.468$ & $.1531$ & $.1555$\\
$(1/\pi,\sqrt2-1)$ & $\frac1{3\pi}=.1061$ & $-.330$ & $-.113$ & $.104$ & $.321$ & $.1067$ & $.1054$\\
$(e-2,\frac{\sqrt5-1}2)$ & $\frac1{3\pi}=.1061$ & $-.314$ & $-.111$ & $.120$ & $.323$ & $.1001$ & $.1018$\\\bottomrule
\end{tabular}
\end{center}
(Here $\beta$ and $r$ are replaced by dyadic approximations with error below $2^{-400}$, which does not affect $D_N$ for $N=10^{30}$ except on a set of
negligible probability.)
\begin{thebibliography}{99}
\bibitem{Bi99} P.~Billingsley, \emph{Convergence of Probability Measures}, 2nd ed., Wiley, 1999.
\bibitem{Bo23a} B.~Borda, Limit laws of maximal Birkhoff sums for circle rotations via quantum modular forms, \emph{Int.\ Math.\ Res.\ Not.\ IMRN} 2023, no.~22, 19340--19389; arXiv:2303.07796.
\bibitem{Bo25} B.~Borda, Equidistribution of continued fraction convergents in $\mathrm{SL}(2,\Z_m)$ with an application to local discrepancy, \emph{J.\ Mod.\ Dyn.} (2025), doi:10.3934/jmd.2025005; arXiv:2303.08504.
\bibitem{DF20} D.~Dolgopyat and B.~Fayad, Deviations of ergodic sums for toral translations II. Boxes, \emph{Publ.\ Math.\ IH\'ES} 132 (2020), 293--352.
\bibitem{DS20} D.~Dolgopyat and O.~Sarig, Quenched and annealed temporal limit theorems for circle rotations, \emph{Ast\'erisque} 415 (2020), 59--85.
\bibitem{E1002} Verifiable Mathematical Reasoning Corpus, Erd\H{o}s Problem \#1002: the Cauchy law for the sawtooth sums $\sum_{k\le n}(\frac12-\{k\alpha\})$, manuscript (2026), included in the same repository as the present paper (directory \texttt{erdos1002}).
\bibitem{HL22} G.~H.~Hardy and J.~E.~Littlewood, Some problems of Diophantine approximation: the lattice-points of a right-angled triangle, \emph{Proc.\ London Math.\ Soc.\ (2)} 20 (1922), 15--36.
\bibitem{IK02} M.~Iosifescu and C.~Kraaikamp, \emph{Metrical Theory of Continued Fractions}, Kluwer, 2002.
\bibitem{Ke60} H.~Kesten, Uniform distribution mod 1, \emph{Ann.\ of Math.\ (2)} 71 (1960), 445--471.
\bibitem{Ke62} H.~Kesten, Uniform distribution mod 1 (II), \emph{Acta Arith.} 7 (1961/62), 355--380.
\bibitem{Kh64} A.~Ya.~Khinchin, \emph{Continued Fractions}, University of Chicago Press, 1964.
\bibitem{Kw26} S.~Kwon, proof claim for Erd\H{o}s Problem \#1002 (manuscript, 2026), with a Lean~4 formalisation by I.~Mian, S.~Siddique and S.~Kwon; S.~Wang, independent proof claim (2026); see \url{https://www.erdosproblems.com/forum/thread/1002/proof-claims}.
\bibitem{Os22} A.~Ostrowski, Bemerkungen zur Theorie der Diophantischen Approximationen, \emph{Abh.\ Math.\ Sem.\ Univ.\ Hamburg} 1 (1922), 77--98.
\bibitem{Sh71} G.~Shimura, \emph{Introduction to the Arithmetic Theory of Automorphic Functions}, Princeton University Press, 1971.
\end{thebibliography}

\end{document}
